\documentclass[12pt]{article}

\usepackage[T1]{fontenc}
\usepackage[utf8]{inputenc}
\usepackage[brazil]{babel}
\usepackage{lmodern}
\usepackage{microtype}
\usepackage[a4paper,margin=2.5cm]{geometry}
\usepackage{setspace}
\usepackage{indentfirst}

\setstretch{1.15}

\title{A Nova Síntese Neoclássica e o Papel da Política Monetária \thanks{Texto Traduzido livremente por Luiz Mario Andrade com o auxílio da ferramenta Chat GPT 5.4}}
\author{Marvin Goodfriend e Robert G. King}
\date{1997}

\begin{document}

\maketitle

\section{1. Introdução}

É comum que a macroeconomia seja retratada como um campo em desordem intelectual, com divergências importantes e persistentes sobre metodologia e substância entre campos rivais de pesquisadores. Uma medida de desordem frequentemente discutida é a distância entre os modelos de preços flexíveis da nova macroeconomia clássica e da análise de ciclos reais de negócios (RBC), na qual a política monetária é essencialmente sem importância para a atividade real, e os modelos de preços rígidos da economia Nova-Keynesiana, na qual a política monetária é vista como central para a evolução da atividade real. Para os formuladores de políticas e os economistas que os assessoram, essa desordem intelectual percebida dificulta o emprego de desenvolvimentos recentes e contínuos na macroeconomia.

As correntes intelectuais dos últimos dez anos são, no entanto, sujeitas a uma interpretação muito diferente: a macroeconomia está se movendo em direção a uma \textit{Nova Síntese Neoclássica}. Na década de 1960, a síntese original envolvia um compromisso com três princípios — por vezes conflitantes: o desejo de fornecer conselhos práticos de política macroeconômica, a crença de que a rigidez de preços no curto prazo estava na raiz das flutuações econômicas e o compromisso de modelar o comportamento macroeconômico usando a mesma abordagem de otimização comumente empregada na microeconomia.

A Nova Síntese Neoclássica herda o espírito da antiga, na medida em que combina elementos keynesianos e clássicos. Metodologicamente, a nova síntese envolve a aplicação sistemática da otimização intertemporal e das expectativas racionais, conforme enfatizado por Robert Lucas. Na síntese, essas ideias são aplicadas às decisões de precificação e produto no coração dos modelos keynesianos, novos e antigos, bem como às decisões de consumo, investimento e oferta de fatores que estão no coração dos modelos clássicos e de RBC. Além disso, a nova síntese também incorpora os \textit{insights} dos monetaristas, como Milton Friedman e Karl Brunner, em relação à teoria e prática da política monetária.

Assim, existem novos fundamentos microeconômicos dinâmicos para a macroeconomia. Essas ideias metodológicas comuns são implementadas em modelos que variam desde os modelos flexíveis e pequenos da pesquisa acadêmica até o novo modelo de política de expectativas racionais do Conselho do Federal Reserve. A Nova Síntese Neoclássica (NNS) sugere um conjunto de conclusões importantes sobre o papel da política monetária. Primeiro, os modelos NNS sugerem que as ações de política monetária podem ter um efeito importante sobre a atividade econômica real, persistindo por vários anos, devido ao ajuste gradual dos preços individuais e do nível geral de preços. Segundo, mesmo em cenários com ajuste de preços dispendioso, os modelos sugerem pouco \textit{trade-off} de longo prazo entre inflação e atividade real. Terceiro, os modelos sugerem ganhos significativos ao eliminar a inflação, decorrentes do aumento da eficiência das transações e da redução das distorções de preços relativos. Quarto, os modelos implicam que a credibilidade desempenha um papel importante na compreensão dos efeitos da política monetária. Essas quatro ideias são consistentes com as declarações públicas de banqueiros centrais de uma ampla gama de países.

Além dos pontos gerais, os modelos NNS permitem a análise de regras alternativas de política monetária dentro de um cenário de expectativas racionais. É nesse papel que eles podem informar — em vez de confirmar — as crenças (\textit{priors}) dos banqueiros centrais. A credibilidade da política monetária parece exigir intuitivamente uma regra simples e transparente. Mas qual delas? Utilizamos a abordagem NNS para desenvolver um conjunto de princípios e diretrizes práticas para uma política monetária \textit{neutra}, definida como aquela que sustenta o produto em seu nível potencial em um ambiente de preços estáveis. A nova síntese sugere que tal política monetária envolve a estabilização do \textit{markup} (margem) médio do preço sobre o custo marginal. Por sua vez, isso implica um regime de política monetária de metas de inflação, que variam relativamente pouco ao longo do tempo. Embora a estabilidade de preços tenha sido sugerida há muito tempo como um objetivo primário para a política monetária, uma série de questões importantes surgiu sobre sua desejabilidade na prática. Confrontamos uma gama de questões de implementação, incluindo a resposta a choques nos preços das commodities, as defasagens longas e variáveis entre a política monetária e o nível de preços, o potencial \textit{trade-off} entre a variabilidade de preços e de produto, e o uso de uma taxa de juros de curto prazo como instrumento de política.

A organização da nossa discussão é a seguinte. Na Seção 2, descrevemos a abordagem geral da síntese neoclássica original conforme articulada por Paul Samuelson. Na Seção 3, revisamos por que a síntese neoclássica original nunca foi totalmente aceita pelos monetaristas, mesmo no auge de sua influência na década de 1960, e depois foi mais fundamentalmente desafiada pela revolução das expectativas racionais. Em seguida, voltamo-nos para trabalhos mais recentes em macroeconomia abrangendo modelos RBC na Seção 4, e a economia Nova-Keynesiana na Seção 5.

A NNS é introduzida e descrita na Seção 6. Analisamos o efeito da política monetária dentro da nova síntese usando duas abordagens complementares. Primeiro, empregamos o método keynesiano padrão que vê a política monetária afetando a demanda agregada real. Segundo, usamos uma alternativa no estilo RBC que vê as variações no \textit{markup} médio como uma fonte de variações na oferta agregada; essas variações de \textit{markup} são análogas aos efeitos dos choques tributários nos modelos RBC. Usamos os \textit{insights} das seções anteriores para desenvolver princípios para a política monetária na Seção 7 e diretrizes práticas para a política monetária na Seção 8. A Seção 9 é um resumo e conclusão.

\section{2. A Síntese Neoclássica}

Conforme popularizada por Paul Samuelson, a síntese neoclássica foi anunciada como um motor de análise que oferecia uma visão keynesiana da determinação da renda nacional — ciclos de negócios surgindo de mudanças na demanda agregada devido à rigidez de salários e preços — e princípios neoclássicos para orientar a análise microeconômica. Em nossa discussão sobre a síntese neoclássica, consideramos três questões principais: a natureza do mecanismo de transmissão monetária, a interação da inflação com a atividade real e o papel da política monetária.

\subsection{2.1 O Mecanismo de Transmissão Monetária}

O arcabouço macroeconômico básico da síntese neoclássica era o modelo IS-LM. A síntese neoclássica gerou uma série de avanços nas décadas de 1950 e 1960 para tornar esse arcabouço mais consistente com a escolha individual e incorporar os elementos dinâmicos que eram tão evidentemente necessários para a modelagem econométrica de séries temporais macroeconômicas. O trabalho teórico racionalizou a demanda por moeda como decorrente da escolha individual na margem, levando a uma explicação microeconômica da taxa de juros e das variáveis de escala no setor monetário. A síntese estimulou avanços na teoria do consumo e do investimento baseados na escolha individual ao longo do tempo. O trabalho econométrico sobre demanda por moeda e investimento desenvolveu especificações de ajuste parcial dinâmico.

Esses novos elementos foram introduzidos em modelos de larga escala da macroeconomia. Nossa discussão foca no modelo MPS do Sistema do Federal Reserve, que foi desenvolvido porque "nenhum modelo existente tem como propósito principal a quantificação da política monetária e seus efeitos na economia", conforme relataram de Leeuw e Gramlich (1968, p. 11). O modelo MPS incluía inicialmente os elementos centrais do arcabouço IS-LM: um bloco financeiro, um bloco de investimento e um bloco de consumo-inventário. A estrutura das possibilidades de produção e a natureza da dinâmica salário-preço eram vistas como importantes, mas secundárias no estágio inicial de desenvolvimento do modelo. Em relação a outros modelos existentes na época, o modelo MPS sugeria efeitos maiores da política monetária porque incorporava um efeito significativo das taxas de juros de longo prazo sobre o investimento e suas defasagens estimadas na demanda por moeda sugeriam um ajuste muito mais rápido do que em modelos anteriores.

Em sua forma totalmente desenvolvida, por volta de 1972, o modelo MPS incorporou várias características estruturais que vale a pena destacar. Foi projetado para ter propriedades de longo prazo como as do modelo de crescimento de consenso de Robert Solow, incluindo a especificação de uma função de produção agregada implicando uma participação constante do trabalho na renda nacional diante do crescimento da produtividade tendencial. Conforme explicado em Ando (1974), no entanto, o modelo MPS tinha uma função de produção de curto prazo que ligava o produto ao insumo trabalho aproximadamente de um para um, como resultado de variações na utilização do capital. A motivação empírica para essa característica é exibida na Figura 1: ao longo dos ciclos de negócios, o total de horas-homem e o produto exibem amplitude semelhante, com a utilização da capacidade medida fortemente pró-cíclica. Na maior parte, essas variações cíclicas nas horas totais decorrem em grande parte de variações no emprego, em vez de horas médias por trabalhador.

\subsection{2.2 Inflação e Atividade Real}

Nos primeiros anos da síntese neoclássica, modelos macroeconométricos foram construídos e análises de política prática foram realizadas assumindo que salários nominais e preços evoluíam independentemente da atividade real e de seus determinantes. De fato, na década de 1950, havia relativamente pouca variabilidade na inflação. Em meados da década de 1960, essa premissa não podia mais ser mantida — a inflação tornou-se uma preocupação política séria e era evidente que a inflação estava relacionada aos desenvolvimentos na economia.

A curva de Phillips tornou-se assim uma parte central da modelagem macroeconômica e da análise de política. Os modelos macroeconômicos eram fechados com setores de salários e preços que indicavam \textit{trade-offs} importantes entre a taxa de inflação e o nível de atividade real. O modelo MPS especificava que o nível de preços era determinado por um \textit{markup} do preço sobre o custo marginal, sendo a taxa de salário nominal um determinante central do custo. Além disso, o modelo MPS fazia o \textit{markup} depender da extensão da utilização e permitia que o nível de preços se ajustasse gradualmente em direção ao custo marginal (Ando, 1974, pp. 544, 552). A versão MPS da curva de Phillips também especificava que a taxa de inflação salarial dependia da taxa de desemprego e da taxa defasada de variação dos preços nominais. Com essas três suposições tomadas em conjunto, como em de Menil e Enzler (1972), o modelo MPS sugeria que o efeito de reduzir a taxa de inflação de longo prazo de 5\% para 0\% era um aumento na taxa de desemprego de 3,5\% para 7\%.

A natureza do \textit{trade-off} entre inflação e desemprego tornou-se central para a política macroeconômica, bem como para a modelagem macroeconômica. Conselheiros de política preocupavam-se com uma espiral salário-preço e temiam que a inflação pudesse desenvolver um ímpeto próprio, como pareceu ser o caso na recessão de 1957-58 (Okun et al., 1969, p. 96; Okun, 1970, p. 8). Pelos padrões de anos posteriores, os resultados para inflação e desemprego foram favoráveis nas décadas de 1950 e 1960. A correlação de Phillips manteve-se notavelmente bem ao longo da década de 1960. No entanto, os assessores econômicos que operavam dentro da tradição da síntese eram pessimistas quanto às perspectivas de domar a inflação.

\subsection{2.3 O Papel da Política Monetária}

Os praticantes da síntese neoclássica viam necessidade de uma gestão ativista da demanda agregada. Dada a escala de rigidez do nível de preços de curto prazo incorporada na síntese neoclássica, a política monetária era reconhecida como tendo efeitos potencialmente poderosos. No entanto, na prática, os conselheiros de política que trabalhavam dentro da síntese viam a política monetária desempenhando um papel permissivo no apoio às iniciativas de política fiscal. Além disso, os economistas consideravam o efeito das taxas de mercado sobre os componentes da demanda agregada sensíveis aos juros como menos importante do que os efeitos diretos do crédito (Okun et al., 1969, pp. 85-92). Eles pensavam que a política monetária funcionava principalmente afetando a disponibilidade de crédito de intermediários financeiros, com importância particular atribuída ao efeito sobre os \textit{spreads} entre as taxas de mercado e as taxas de depósito então regulamentadas. Consequentemente, havia uma relutância em deixar o fardo da política de estabilização recair sobre a política monetária, uma vez que ela funcionava por meio de uma espécie de distorção.

Apesar da relutância em usá-la, os praticantes da síntese neoclássica reconheciam que a política monetária poderia controlar a inflação. A visão de Okun (1970, p. 8) era representativa: "a cura básica para a inflação é remover ou compensar sua causa: cortar a demanda agregada por meio de política fiscal ou monetária suficientemente para que os gastos monetários não excedam mais o valor dos bens". James Tobin pôde dizer sobre o aperto da política monetária de 1966 para combater a inflação em rápida ascensão que "o fardo da contenção recaiu quase inteiramente sobre o Fed, que agiu vigorosa e corajosamente".

Assim, a política monetária na síntese neoclássica era considerada um instrumento poderoso, mas mal adaptado para controlar a inflação ou realizar política de estabilização. Embora a política monetária pudesse controlar a inflação em teoria, a visão prática era que a inflação era governada principalmente por fatores psicológicos e inércia, de modo que a política monetária poderia ter apenas um efeito muito gradual. Como a política monetária criava distorções entre setores, a política fiscal era mais adequada para controlar o ciclo de negócios.

\section{3. Monetarismo e Expectativas Racionais}

Quando surgiu na década de 1960, o monetarismo pareceu ameaçar a síntese neoclássica. Em parte, isso ocorreu porque os monetaristas se retratavam como descendentes intelectuais da teoria quantitativa da moeda pré-keynesiana, conforme articulada por Irving Fisher e outros. Em parte, foi porque os monetaristas questionaram grande parte da doutrina da síntese, por exemplo, a eficácia da política fiscal e a estabilidade estrutural da curva de Phillips. Nas décadas de 1970 e 1980, muitos \textit{insights} monetaristas seriam incorporados à síntese de base ampla, e por uma boa razão: o monetarismo era um conjunto de princípios para conselhos práticos de política, estava comprometido com o raciocínio neoclássico e também identificava a fonte dos ciclos de negócios na rigidez do nível de preços no curto prazo. No entanto, ao mesmo tempo, a crítica de Lucas à política macroeconométrica e a subsequente introdução de expectativas racionais na macroeconomia levaram a um questionamento mais amplo da síntese neoclássica.

A teoria quantitativa — o coração do monetarismo — sugeria organizar a análise monetária em termos da oferta de moeda nominal e da demanda por saldos monetários reais. Esse foco tinha implicações para o mecanismo de transmissão monetária, para a ligação entre inflação e atividade real, e para o papel da política monetária.

\subsection{3.1 O Mecanismo de Transmissão Monetária}

O arcabouço monetarista básico era a equação quantitativa, que introduzimos usando a notação que manteremos ao longo do artigo. De acordo com a teoria quantitativa, a renda nominal ($Y_t$) é o resultado do estoque de moeda ($M_t$) e sua velocidade ($v_t$):

\begin{equation}
\log Y_t = \log M_t + \log v_t. \tag{3.1}
\end{equation}

Os monetaristas tornaram a teoria quantitativa operacional ao considerar a moeda como autônoma. Eles também construíram um modelo econométrico com base em seu arcabouço analítico. O modelo de St. Louis de Anderson e Jordan (1968) era simplesmente a equação quantitativa em um contexto de defasagem distribuída, com uma especificação flexível introduzida para capturar o ajuste dinâmico da demanda e oferta de moeda.

A visão monetarista do mecanismo de transmissão estava nitidamente em desacordo com a síntese neoclássica, que tendia a ver os principais canais de transmissão funcionando através da disponibilidade de crédito e, em segundo lugar, através do efeito das taxas de juros de longo prazo sobre o investimento. Os monetaristas consideravam ambos os canais como secundários. Eles focavam na moeda em vez dos canais de crédito.

Seguindo Irving Fisher, os monetaristas reconheciam que as taxas de juros nominais continham um componente real e um prêmio pela inflação esperada. Como outras defasagens, aquelas na formação de expectativas eram consideradas longas e variáveis. Como questão prática, no entanto, os monetaristas consideravam a maior parte da variação nas taxas de longo prazo como reflexo de prêmios de inflação, conferindo às taxas de longo prazo um papel relativamente menor na transmissão da política monetária para a atividade real.

\subsection{3.2 Inflação e Atividade Real}

Os monetaristas também divergiam em sua visão sobre a ligação entre inflação e atividade real. Na maior parte, os monetaristas reconheciam que não possuíam uma teoria confiável para prever a divisão de curto prazo do crescimento da renda nominal entre o nível de preços ($P_t$) e o produto real ($y_t$) — eles não tinham uma equação de preços de curto prazo. De várias maneiras, eles interpretaram a aparente não neutralidade da moeda no curto prazo como resultado da rigidez do nível de preços. Mas observaram que o efeito das ações de política monetária na economia era longo e variável. Eles tendiam a atribuir essa variabilidade a diferenças no grau em que as ações de política eram esperadas, porque as expectativas determinavam o grau em que preços e salários se ajustariam para neutralizar uma injeção de moeda.

Essas considerações sobre expectativas foram explicitadas por Friedman (1968), que descreveu como o ajuste incompleto das expectativas poderia levar salários e preços a responder lentamente a mudanças na moeda. Ao mesmo tempo, Friedman sugeriu que a inflação sustentada não deveria afetar a atividade real no longo prazo, definido como uma situação na qual as expectativas estivessem corretas, uma vez que o produto seria então determinado por forças reais. As sugestões de Friedman foram oportunas. Como mostrado na Figura 2, a inflação aumentou drasticamente na década de 1970 com pouca expansão acompanhante da atividade real.

\subsection{3.3 O Papel da Política Monetária}

Os monetaristas também viam um papel dramaticamente diferente para a política monetária. Desconfiados de uma política monetária discricionária e ativista, eles buscaram formular regras fixas simples para a política. Com a interpretação de Friedman e Schwartz (1963) da Grande Depressão em mente, eles acreditavam que a autoridade monetária deveria evitar grandes choques monetários na macroeconomia, sugerindo uma regra na qual a quantidade de moeda crescesse a uma taxa constante suficiente para acomodar o crescimento da produtividade tendencial (Friedman, 1960). Após argumentar que a inflação sustentada tem pouco efeito sobre a atividade real, Friedman (1969) descreveu um regime monetário de longo prazo que envolvia deflação sustentada, tornando a taxa de juros nominal zero e, assim, proporcionando uma quantidade ótima de moeda.

Na prática, também houve diferenças importantes no papel sugerido para a política monetária ao longo do ciclo de negócios. Enquanto os assessores de política da síntese neoclássica buscavam que o Federal Reserve mantivesse as taxas de juros inalteradas enquanto a política fiscal era variada, os monetaristas pensavam que a suavização da taxa de juros contribuía para flutuações na atividade econômica real ao fazer o estoque de moeda variar pró-ciclicamente.

\subsection{3.4 Expectativas Racionais}

Como foi o caso do monetarismo, a introdução de expectativas racionais na macroeconomia no início da década de 1970 a princípio pareceu incompatível com a síntese neoclássica. Isso foi particularmente irônico, pois John Muth motivou sua hipótese de expectativas racionais sugerindo que os indivíduos formam expectativas de maneira ótima, o que é uma extensão natural do princípio neoclássico de que a economia é habitada por agentes racionais e maximizadores.

Os primeiros modelos novos clássicos, como o de Sargent e Wallace (1975), incorporaram a visão de Friedman de que variações percebidas na moeda levavam simplesmente a mudanças nos preços, com apenas mudanças monetárias mal percebidas tendo efeitos reais. Combinado com as expectativas racionais, esse forte mecanismo de neutralidade levou a declarações muito específicas e controversas sobre o papel da política monetária. Primeiro, como na análise monetarista, o banco central deveria evitar criar choques monetários. Segundo, uma ampla classe de regras monetárias levava às mesmas flutuações na atividade real, uma vez que os efeitos reais de variações percebidas na moeda seriam neutralizados por movimentos no nível de preços.

\subsection{3.5 Credibilidade}

Mesmo que o resultado da ineficácia da política fosse frágil, outras implicações de longo alcance estendem-se à maioria dos modelos macroeconômicos modernos, incluindo o arcabouço de preços rígidos que discutimos abaixo. O raciocínio das expectativas racionais ensina que o efeito de um dado choque não pode ser calculado sem compreender sua persistência ou a extensão em que ele era esperado e preparado antecipadamente. Esse ponto, entregue com força em Lucas (1976), revolucionou a análise de política, implicando que não se pode prever o efeito de uma ação de política em um ponto no tempo sem levar em conta a natureza do regime de política do qual ela provém.

Sargent (1986) vinculou essas ideias explicitamente à natureza do processo inflacionário: "a inflação apenas \textit{parece} ter um ímpeto próprio. É, na verdade, a política governamental de longo prazo de persistir em grandes déficits e criar moeda a taxas elevadas que confere ímpeto à taxa de inflação". Revisando uma série de episódios históricos nos quais os países tentaram reduzir altas taxas de inflação, ele argumentou que os custos da desinflação — perda de produto — eram muito menores se o compromisso do governo com a desinflação fosse crível do que se não fosse. No entanto, ironicamente, o novo modelo macroeconômico clássico atribuía pouca importância à credibilidade. Nesse modelo, as intenções futuras do banco central são muito importantes para a evolução do nível de preços, porque afetam a inflação esperada, mas são de relevância limitada para a atividade real desde que sejam percebidas com precisão. Consequentemente, embora muitos bancos centrais vissem a credibilidade como importante, eles relutavam em usar o novo modelo macroeconômico clássico para análise de questões de política monetária.

\section{4. Ciclos Reais de Negócios}

Embora as expectativas racionais tenham sido introduzidas na macroeconomia para estudar as ligações entre variáveis reais e nominais, suas implicações foram trabalhadas de forma mais sistemática dentro do programa de pesquisa de ciclos reais de negócios (RBC). A forte neutralidade monetária incorporada nos modelos RBC impediu seu uso generalizado na análise de política macroeconômica até o momento. Mas vemos a lógica RBC como uma parte central da Nova Síntese Neoclássica. Uma razão é que o programa RBC constrói modelos nos quais as políticas alternativas podem ser comparadas com base em medidas de benefícios ou custos de utilidade, em vez de objetivos \textit{ad hoc}. Outra é que o arcabouço RBC permite a análise de choques de política e outros em um contexto dinâmico-estocástico de um sistema totalmente especificado, conforme exigido pelo raciocínio de expectativas racionais. O programa RBC integra e esclarece a substituição intertemporal que está no coração da macroeconomia — envolvendo consumo, investimento e comportamento de oferta de trabalho — e, ao fazê-lo, esclarece os determinantes da taxa de juros real. Finalmente, os modelos RBC fornecem \textit{insights} sobre a natureza das não neutralidades cíclicas nos modelos NNS e também descrevem resultados macroeconômicos sob uma política monetária neutra.

\subsection{4.1 Elementos Centrais dos Modelos RBC}

A abordagem RBC utiliza modelos de equilíbrio geral real para estudar fenômenos macroeconômicos. Um elemento chave é a abordagem de otimização intertemporal para o consumo e a oferta de trabalho. Outro é a análise intertemporal semelhante de investimento e demanda por trabalho, decorrente das decisões de maximização de lucro das firmas. Os planos das famílias e das firmas são então combinados em um equilíbrio geral, no qual quantidades e preços são determinados simultaneamente.

\subsection{4.2 Choques de Produtividade}

O programa RBC focou os macroeconomistas na pró-ciclicidade da produtividade medida dos insumos de fatores. Nas mãos de Prescott (1986) e Plosser (1989), o modelo RBC básico foi visto como capaz de gerar ciclos de negócios que se assemelhavam aos dos EUA e de outras economias quando impulsionado por resíduos de Solow. Para fins de definição desses resíduos e para discussão de outras questões abaixo, escrevemos a função de produção como retornos constantes de escala em trabalho ($n$) e capital ($k$), deslocando-se ao longo do tempo como resultado de choques de produtividade ($a$):

\begin{equation}
y_t = a_t F(n_t, k_t). \tag{4.1}
\end{equation}

No modelo RBC, os choques de produtividade têm dois conjuntos de efeitos sobre o produto. Um é que eles aumentam ou diminuem mecanicamente o produto, como enfatizado por Solow em sua famosa decomposição:

\begin{equation}
\frac{dy_t}{y_t} = s_n \frac{dn_t}{n_t} + s_k \frac{dk_t}{k_t} + \frac{da_t}{a_t}, \tag{4.2}
\end{equation}

onde $s_k$ e $s_n$ são as participações dos fatores trabalho e capital. No entanto, os choques de produtividade também exercem efeitos sobre a atividade macroeconômica, porque afetam as escalas de produto marginal (demanda por fatores). Essas influências marginais (substituição) interagem com a motivação de suavização incorporada nas preferências das famílias para governar a resposta dinâmica da economia. Um aumento temporário na produtividade atual, por exemplo, torna mais valioso para as famílias trabalhar (reduzir o lazer) e investir (adiar o consumo atual). Dentro do modelo RBC, esses mecanismos explicam, por exemplo, a pró-ciclicidade do insumo trabalho e a resposta de alta amplitude do investimento. A abordagem RBC força o pesquisador a explicar a resposta da macroeconomia em termos de efeitos substituição e riqueza sobre as famílias.

Uma questão importante sobre a abordagem RBC tem sido a medição dos choques de produtividade, particularmente se o método de Solow mede incorretamente os insumos de fatores. Pesquisas subsequentes focaram na utilização variável do capital como uma fonte de medição incorreta: trabalhos recentes de Burnside, Eichenbaum e Rebelo (1995) reduzem a variabilidade do resíduo de Solow de forma tão substancial que um adepto da abordagem RBC pode se preocupar com o fato de restar pouco em termos de choques de produtividade.

\subsection{4.3 Racionalizando Altas Elasticidades de Oferta}

Ao focar a atenção no lado da oferta, os modeladores RBC provocaram muitas questões, sendo uma das mais básicas: as variações de oferta de trabalho de alta amplitude assumidas nos modelos RBC são contrafatuais? Os primeiros modelos RBC assumiam que a oferta de trabalho agregada variava apenas por um trabalhador individual (o agente representativo) alterando o número de horas trabalhadas. Esse mecanismo é reconhecidamente inconsistente com a evidência microeconômica sobre a oferta de trabalho.

No entanto, ao longo do ciclo de negócios, ocorrem grandes mudanças no esforço de trabalho. Como ilustrado pela comparação dos painéis (a) e (c) da Figura 1, estas decorrem principalmente de mudanças no número de indivíduos empregados, em vez de mudanças no número de horas trabalhadas por cada indivíduo. Modificações importantes do arcabouço RBC básico modelaram tais movimentos de entrada e saída da força de trabalho, gerando elasticidades de oferta de trabalho agregadas extremamente altas, mantendo pequenas elasticidades micro. Outros estudos recentes apresentam utilização variável do capital, com uma oferta de serviços de capital altamente sensível a mudanças nos preços dos fatores, de modo que a utilização é fortemente pró-cíclica. No geral, a abordagem RBC moderna descreve uma macroeconomia altamente sensível a choques reais. Hall (1991) aponta que muitas abordagens para os ciclos de negócios exigem uma economia de "alta substituição" como aquela construída pelos pesquisadores RBC.

\subsection{4.4 Moeda em Modelos RBC}

No início do programa de pesquisa RBC, um setor monetário foi adicionado para explorar os tipos de correlações de ciclo de negócios entre moeda e produto que poderiam surgir se os choques de produtividade fossem o principal fator impulsionador (King e Plosser, 1984). Em um estágio posterior da pesquisa, os efeitos do imposto inflacionário foram explorados (Cooley e Hansen, 1989). A partir dessa pesquisa e de outros trabalhos ao longo da última década, surgiu uma série de conclusões amplamente compartilhadas pelos macroeconomistas. Primeiro, variações endógenas na oferta de moeda decorrentes das ações conjuntas dos bancos privados e da autoridade monetária explicam, pelo menos em parte, a correlação de ciclo de negócios entre moeda e produto. Segundo, embora versões de modelos RBC suplementadas com um setor monetário possam, em princípio, explicar a correlação entre moeda e produto, elas se saem menos bem ao explicar a variação cíclica nas taxas de juros reais e nominais (Sims, 1992), sugerindo que há mais no ciclo do que choques de produtividade reais que causam variações simpáticas na moeda. Terceiro, as consequências previstas das variações cíclicas na inflação esperada são quantitativamente pequenas dentro de modelos de preços flexíveis, se a demanda por moeda for modelada via \textit{cash-in-advance} ou com uma tecnologia de transações explícita. Ou seja, para fins de ciclo de negócios, um modelo RBC com um mecanismo monetário explícito funciona muito como um modelo RBC com uma função de demanda por moeda apenas anexada após uma análise de equilíbrio geral real.

\subsection{4.5 Análise da Inflação Sustentada}

Estudos sobre os custos da inflação estável realizados sob a rubrica RBC levaram a uma compreensão revisada dos benefícios que podem ser obtidos com a redução da inflação. Uma referência básica nesta área é Lucas (1993), que calcula que o custo de bem-estar de uma inflação de 7\% pode ser de cerca de 1\% do produto usando uma variante do modelo de \textit{shopping-time} da demanda por moeda. Como a tecnologia de transações de Lucas não possui nível de saciedade para saldos monetários, a maior parte de seus ganhos estimados com a redução da inflação para o nível de Friedman (1969) surge como resultado da deflação. No entanto, estimando os parâmetros de um modelo de \textit{shopping-time} com dados anuais dos EUA de 1915 a 1992, Wolman (1996) conclui que a experiência dos EUA parece mais consistente com uma tecnologia de transações com um nível de saciedade para saldos monetários. Este modelo alternativo de demanda por moeda fornece aproximadamente o mesmo ganho total ao reduzir a inflação, mas localiza a maior parte dele entre 7\% e zero de inflação.

\subsection{4.6 Política Fiscal e Choques Fiscais em um Cenário RBC}

Outro tópico importante da análise RBC tem sido o estudo da política fiscal e dos distúrbios fiscais em equilíbrio geral real. No modelo RBC, mudanças nas alíquotas de impostos têm um efeito poderoso sobre a atividade real. Em particular, variações em um imposto abrangente sobre a renda ou vendas afetam os retornos reais dos fatores trabalho e capital após impostos, induzindo substituições entre bens e ao longo do tempo que influenciam as quantidades de esforço de trabalho e investimento escolhidas por um agente representativo. Por exemplo, o salário real após impostos é:

\begin{equation}
w_t = (1 - \tau_t) a_t \frac{\partial F(n_t, k_t)}{\partial n_t}, \tag{4.3}
\end{equation}

onde $\tau_t$ é a taxa de imposto na data $t$ e $w_t$ é o salário real na data $t$. Assim, do ponto de vista do retorno marginal ao trabalho, o imposto funciona exatamente como um choque de produtividade. Consequentemente, mudanças nos impostos abrangentes sobre a renda exercem uma influência de alta octanagem no modelo RBC.

Estudos RBC de choques fiscais reais dos EUA, como o de McGrattan (1994), chegam a uma conclusão irônica. Mudanças nas alíquotas de impostos têm efeitos poderosos sobre a atividade macroeconômica, mas como a variação nas taxas de impostos medidas sobre o capital e a renda nos EUA em frequências de ciclo de negócios é pequena, esses choques não contribuem muito para a variabilidade geral do ciclo de negócios. No entanto, veremos abaixo que as mudanças nos \textit{markups} podem ser interpretadas como impostos de uma forma potencialmente volátil ciclicamente.

\section{5. Economia Nova-Keynesiana}

A abordagem Nova-Keynesiana para a macroeconomia evoluiu em resposta à controvérsia monetarista e a questões fundamentais levantadas pela crítica de Lucas, e a fim de fornecer uma alternativa ao arcabouço competitivo de preços flexíveis da análise RBC. Nossa discussão deste amplo programa de pesquisa será dividida em três partes. Primeiro, revisamos os trabalhos iniciais de Gordon (1982) e Taylor (1980). Em seguida, discutimos os fundamentos microeconômicos Nova-Keynesianos mais recentes, que destacam a competição monopolística e o ajuste de preços custoso. Finalmente, focamos no ajuste de preços otimizador em um cenário dinâmico.

\subsection{5.1 Modelos Nova-Keynesianos de Primeira Geração}

Nos modelos Nova-Keynesianos de primeira geração, Gordon (1982) e Taylor (1980) modernizaram a especificação do bloco salário-preço para incorporar \textit{insights} monetaristas e de expectativas racionais.

\subsubsection{5.1.1 A Equação de Preços de Gordon} No lado empírico, Gordon (1982) estimou a dinâmica de preços usando uma exogeneidade aproximada monetarista da demanda agregada nominal. Abstendo-se da consideração separada dos salários nominais porque via sua dinâmica como essencialmente idêntica à dos preços, Gordon estimou equações de preços da forma:

\begin{equation}
\pi_t = \lambda(L)\pi_{t-1} + G(\log Y_t - \log Y_{t-1}) + ps_t + \eta_t, \tag{5.1}
\end{equation}

onde $\pi_t = \log P_t - \log P_{t-1}$ é a taxa de inflação, $\lambda(L)$ é um polinômio no operador de defasagem, $\log Y_t - \log Y_{t-1}$ é o crescimento da renda nominal, $ps_t$ captura os efeitos de choques de preços observáveis, e $\eta_t$ é um termo de erro.

Gordon interpretou os coeficientes $\lambda(L)$ como indicadores de como o nível de preços se ajusta gradualmente em direção a um nível de longo prazo exigido pela renda nominal e um nível de atividade real de "taxa natural". Houve três descobertas principais da investigação de Gordon: Primeiro, houve um valor de $G$ numericamente pequeno na equação de preços. Estimando equações de preços trimestrais ao longo de quase um século de dados e várias subamostras, Gordon encontrou coeficientes de inclinação na faixa de $G = 0,10$, indicando um pequeno efeito de impacto do produto sobre os preços igual a $G/(1+G) = 0,09$. Segundo, as defasagens foram estimadas como sendo muito importantes na equação de preços: a defasagem média entre produto e preços foi de mais de um ano. Gordon interpretou isso como evidência de ajuste gradual do nível de preços a mudanças nos gastos nominais.

No entanto, Gordon também encontrou mudanças notáveis em suas estimativas quando os noventa anos de dados foram divididos em três ou mais subperíodos. Na subamostra inicial que vai de 1892 a 1929, houve mudanças importantes nos efeitos da renda nominal durante o período de guerra 1915-1922. Em particular, o coeficiente estimado sobre a renda nominal aumentou substancialmente, com uma grande diferença surgindo entre o crescimento da renda nominal esperado ($G = 0,47$) e o crescimento da renda nominal inesperado ($G = 0,25$). Medidas de choques de oferta — notadamente preços de energia e commodities — tornaram-se cada vez mais importantes no período da amostra pós-Segunda Guerra Mundial. Finalmente, a soma dos coeficientes sobre a inflação defasada, $\lambda(1)$, aumentou substancialmente de 0,4 durante 1892-1929 para mais de 1 durante 1954-1980.

\subsubsection{5.1.2 Abordagem de Expectativas Racionais de Taylor para a Fixação de Salários} O mais resistente dos modelos macroeconômicos de expectativas racionais Nova-Keynesianos de primeira geração é o de Taylor (1980). Na terminologia moderna, a visão de Taylor era que a firma e seus trabalhadores estabeleciam um salário fixo ao longo da vida de um contrato de período $J$. Assumia-se que as negociações salariais eram escalonadas ao longo do tempo com $1/J$ dos contratos estabelecidos a cada período. A representação matemática mais simples do mecanismo de fixação de salários de Taylor é a seguinte. O salário nominal fixado na data $t$, $\log W_t^*$, depende do nível médio de preços esperado ao longo do contrato, $(1/J) \sum_{j=0}^{J-1} E_t \log P_{t+j}$; da rigidez média do mercado de trabalho [incorporada como $(h/J) \sum_{j=0}^{J-1} E_t e_{t+j}$, onde $e_t$ é a rigidez do mercado de trabalho na data $t$ e $h$ governa a resposta salarial a essa rigidez]; e de um choque salarial ($\nu_t$):

\begin{equation}
\log W_t^* = \frac{1}{J} \sum_{j=0}^{J-1} E_t \log P_{t+j} + \frac{h}{J} \sum_{j=0}^{J-1} E_t e_{t+j} + \nu_t. \tag{5.2}
\end{equation}

Taylor (1980) adotou um modelo macroeconômico muito simples para focar nas consequências desse comportamento de fixação de salários. Primeiro, Taylor especificou que o nível de preços era uma média simples de salários, motivada pela referência a um monopolista com custo marginal constante selecionando um \textit{markup} fixo,

\begin{equation}
\log P_t = \frac{1}{J} \sum_{j=0}^{J-1} \log W_{t-j}^*. \tag{5.3}
\end{equation}

Segundo, assim como Gordon (1982), Taylor fez a suposição monetarista de que os gastos nominais eram determinados por uma equação quantitativa. Terceiro, Taylor assumiu que a rigidez do mercado de trabalho estava relacionada ao produto: $e_t = g_1 \log y_t$. Quarto, Taylor assumiu uma regra de estoque de moeda ativista para a política monetária, especificamente que

\begin{equation}
\log M_t = g_2 \log P_t. \tag{5.4}
\end{equation}

Tal como aconteceu com os modelos de expectativas racionais novos clássicos anteriores, o modelo de expectativas racionais de Taylor exigia a especificação do comportamento da autoridade monetária. Em vez de considerar a autoridade monetária como uma fonte de impulsos de ciclo de negócios, ele a viu como ajustando o estoque de moeda ao nível de preços com um coeficiente de resposta $g_2$.

\subsubsection{5.1.3 Implicações de Ciclo de Negócios e Política do Arcabouço de Taylor} Existem quatro implicações do arcabouço de Taylor. Primeiro, Taylor produziu um padrão em "formato de corcova" da dinâmica cíclica do produto (desemprego) em resposta a choques salariais $\nu$, o que Taylor sugeriu ser uma medida de sucesso, porque vários pesquisadores empíricos haviam estimado modelos de séries temporais que implicavam tais perfis. Segundo, Taylor demonstrou que a regra de política importava para a evolução da atividade real. Terceiro, Taylor destacou um novo \textit{trade-off} de política monetária entre a variabilidade do produto e a variabilidade da inflação dentro de seu modelo, mesmo com a suposição mantida de que não havia \textit{trade-off} de longo prazo entre a taxa de inflação e o nível do produto. Se os choques de velocidade fossem pequenos, por exemplo, então um banco central poderia eliminar amplamente a variabilidade real acomodando os movimentos do nível de preços ($g_2$ próximo de um), mas isso exigiria maior variabilidade no nível de preços. Quarto, ele mostrou que as expectativas racionais importavam muito — para a resposta da economia a choques e para o desenho de regras de política monetária — ao contrastar seus resultados com aqueles baseados em expectativas extrapolativas.

Importante notar, modelos de expectativas racionais de salários rígidos e preços rígidos como o de Taylor também explicaram as principais descobertas de Gordon, pelo menos em forma ampla. Defasagens de salários nominais e preços eram variáveis de estado importantes nesses modelos, refletindo o ajuste gradual a choques reais e nominais. Além disso, os efeitos de variações aproximadamente exógenas na renda nominal dependiam de maneira central de quão persistentes se esperava que fossem, uma vez que (5.2) indicava que as expectativas de preços desempenhavam um papel fundamental na fixação de salários.

Tomado em conjunto com seu trabalho subsequente em modelos macroeconométricos maiores incorporando o ajuste gradual de preços, o modelo teórico de Taylor teve um grande impacto intelectual. No entanto, ao mesmo tempo, havia um mal-estar sobre os modelos de salários escalonados de Taylor. Nos Estados Unidos, em particular, apenas uma pequena parcela da força de trabalho estava sujeita a contratos explícitos de múltiplos períodos. Além disso, os fundamentos microeconômicos do processo de fixação de salários eram superficiais.

\subsection{5.2 Modelos Nova-Keynesianos de Segunda Geração}

No próximo estágio da pesquisa, os economistas Nova-Keynesianos deslocaram a localização da rigidez nominal dos salários para os preços. Nesse novo trabalho, as firmas fixadoras de preços foram explicitamente modeladas como competidores monopolísticos. O arcabouço de competição imperfeita foi usado para explicar o efeito real do produto da moeda quando os preços estavam sujeitos a custos de ajuste, para desenvolver vários mecanismos de amplificação e para destacar os potenciais custos sociais dos ciclos de negócios.

\subsubsection{5.2.1 Modelos de Competição Monopolística Explícita} Durante a década de 1980, as implicações da competição monopolística foram exploradas em uma ampla gama de campos, incluindo crescimento econômico, comércio internacional e finanças, e macroeconomia. Em cada caso, a competição imperfeita prometia a compreensão de questões que eram intrigantes sob a perspectiva da teoria competitiva. Na macroeconomia, a competição monopolística era importante para analisar como as firmas estabelecem preços. No cenário competitivo padrão, as firmas aceitam os preços de mercado como dados e ajustam a quantidade em resposta a variações de preços e custos. Em contraste, em Blanchard e Kiyotaki (1987) e Rotemberg (1987), as firmas são competidores monopolísticos e estabelecem preços a fim de maximizar o lucro. Esses estudos consideram o consumo como um agregado de um contínuo de produtos diferenciados, $c_t = [\int_0^1 c_t(z)^{(\varepsilon-1)/\varepsilon} dz]^{\varepsilon/(\varepsilon-1)}$. Uma firma individual produzindo o produto $z$ enfrenta a demanda de elasticidade constante

\begin{equation}
c_t(z) = c_t \left( \frac{P_t(z)}{P_t} \right)^{-\varepsilon}, \tag{5.5}
\end{equation}

que é deslocada pelo nível de preços agregado e pelo nível de demanda de consumo agregada. Investimento e compras governamentais poderiam ser vistos de forma semelhante como agregados de produtos diferenciados, levando a uma versão de (5.5) que substituiu $c_t$ por uma medida de demanda agregada. A forma implícita do índice de preços (perfeito) associado aos gastos agregados é

\begin{equation}
P_t = \left( \int_0^1 P_t(z)^{1-\varepsilon} dz \right)^{1/(1-\varepsilon)}. \tag{5.6}
\end{equation}

Consequentemente, com um custo marginal nominal de $\Psi_t$, uma firma otimizadora estabeleceria seu preço a um \textit{markup} constante sobre o custo marginal, $P_t(z) = [\varepsilon/(\varepsilon-1)] \Psi_t$, sendo o \textit{markup} dado pela fórmula convencional. Assim, a competição monopolística racionaliza uma firma estabelecendo um preço e estabelecendo-o em um nível maior que o custo marginal. A competição imperfeita não racionaliza, por si só, a rigidez nominal.

\subsubsection{5.2.2 Incorporação da Rigidez Nominal} No nível microeconômico, a rigidez dos preços nominais é uma característica de nossa vida cotidiana. Assim, se estamos desenvolvendo "fundamentos micro para a macroeconomia", é importante ter modelos que possam explicar essas práticas de precificação observadas. A explicação mais direta é que pequenos custos reais de mudança de preços nominais — custos de menu — explicam os preços rígidos. É uma questão aberta se pequenos custos de menu podem levar a uma rigidez sustentada dos preços de bens individuais, particularmente em uma situação de inflação positiva. Na maior parte, na abordagem de modelagem Nova-Keynesiana, o ajuste discreto e ocasional de preços individuais é simplesmente uma característica do ambiente, racionalizada de formas mais ou menos elaboradas. Neste artigo, como nessa literatura, focamos menos em por que os preços individuais podem ser estabelecidos antecipadamente e mais nas implicações que o ajuste de preços individual discreto e ocasional tem para o comportamento do nível de preços agregado e da atividade econômica real.

\subsubsection{5.2.3 Causas e Consequências dos Ciclos de Negócios Monetários} Os economistas Nova-Keynesianos também enfatizaram que a competição imperfeita é importante para o efeito da moeda sobre o produto se houver rigidez de preços nominais. Para ver o poder deste argumento, pense no caso da competição perfeita. Se a demanda aumenta, mas o preço permanece o mesmo, a firma não responderá, encaminhando seus potenciais clientes para outro lugar. Em contraste, se seu preço for fixado em um nível que exceda o custo marginal, então é desejável que uma firma individual expanda seu produto se sua demanda aumentar. O caso fácil é se o custo marginal não estiver relacionado ao produto da firma, pois então ela absorverá toda a variação da demanda sem sofrer um declínio em seu \textit{markup}. Mesmo que o custo marginal aumente com o produto, seja ao nível da firma ou em equilíbrio geral, continuará sendo lucrativo satisfazer a demanda enquanto o preço exceder o custo marginal. Em resposta a uma expansão econômica geral — um aumento em $c_t$ em (5.5) acima — é plausível que o custo marginal aumente porque as firmas devem pagar salários reais mais altos para garantir o insumo trabalho para produzir produto adicional. Consequentemente, os Novos Keynesianos destacam a importância de movimentos pró-cíclicos nos salários reais e no custo marginal.

Como questão relacionada, a análise Nova-Keynesiana também sugere um novo conjunto de conclusões para a análise de bem-estar do ciclo de negócios. Com competição monopolística, o poder de mercado das firmas significa que há um nível muito baixo de emprego e produto em média. A análise Nova-Keynesiana fornece, assim, um relato coerente da tentação de expandir a economia presente na literatura sobre políticas monetárias inconsistentes no tempo (Barro e Gordon, 1983). Além disso, os formuladores de políticas monetárias não devem ser indiferentes a mudanças de curto prazo no emprego que surgem de mudanças na moeda quando os preços são rígidos. Notavelmente, uma diminuição no emprego e no produto que resulte de uma política monetária contracionista reduz o bem-estar do indivíduo representativo ao aumentar as distorções monopolísticas.

\subsubsection{5.2.4 Origens e Implicações da Competição Monopolística} Existe uma gama de mecanismos econômicos, é claro, que são consistentes com a competição monopolística. Para nós, o mais plausível é que as firmas enfrentem custos fixos importantes, incluindo custos gerais de administração (\textit{overhead}). Isso sugere a modificação da função de produção para

\begin{equation}
y_t = a_t [F(n_t, k_t) - \Phi], \tag{5.7}
\end{equation}

onde $\Phi$ é uma medida de custos fixos, que plausivelmente se assume exibirem os mesmos requisitos de intensidade de fatores e mudanças técnicas que governam o produto final. Com tal função de produção, a firma representativa tem custo marginal constante (a preços de fatores dados) e custo médio decrescente.

Hall (1988) demonstra que a decomposição de Solow modificada é

\begin{equation}
\frac{dy_t}{y_t} = (1 + \phi) \left( s_n \frac{dn_t}{n_t} + s_k \frac{dk_t}{k_t} \right) + \frac{da_t}{a_t}, \tag{5.8}
\end{equation}

onde $s_n$ e $s_k$ são participações no custo total e $\phi$ é a razão entre custo fixo e custo variável. Essa decomposição destaca as consequências dos custos fixos. Primeiro, o resíduo de Solow padrão varia com o ciclo de negócios, mesmo que não haja choques de produtividade. Segundo, há um mecanismo de amplificação, de modo que uma mudança de um por cento no trabalho altera o produto em $(1 + \phi)s_n$ por cento.

A abordagem Nova-Keynesiana permite uma ampla gama de suposições sobre a natureza e extensão da competição imperfeita. Se não houver lucros de monopólio puros, então o \textit{markup} do preço sobre o custo marginal deve simplesmente cobrir os custos fixos, ou seja, devemos ter $\mu = 1 + \phi$ em média, o que assumimos ao longo de nossa discussão. Em vários exercícios quantitativos abaixo, também será necessário assumirmos uma posição sobre o valor do \textit{markup} de estado estacionário. Comparado a alguns outros estudos recentes, adotamos um valor pequeno, $\mu = 1,1$, o que corresponde a um \textit{markup} "líquido" de 10\% e uma elasticidade de demanda de cerca de 11. Fazemos isso por duas razões. Primeiro, é amplamente consistente com os \textit{markups} observados nas indústrias de construção e serviços automotivos, ou seja, \textit{markups} na faixa de 7\% a 15\% em contratos e faturas de venda. Segundo, é consistente com os estudos empíricos detalhados de Basu e Fernald (1997).

\subsection{5.3 Modelos Dinâmicos de Precificação}

Modelos de dinâmica de preços baseados em custos reais fixos de mudança de preços nominais foram desenvolvidos pela primeira vez no início da década de 1970. Nesses modelos, as firmas escolhem o momento e a magnitude de seus ajustes de preços em resposta ao estado da economia, incluindo a taxa média de inflação e o estágio do ciclo de negócios. Essa abordagem de precificação dependente do estado (\textit{state-dependent}) é atraente do ponto de vista microeconômico porque (1) observa-se que as firmas individuais ajustam discretamente seus preços em intervalos infrequentes de duração aparentemente estocástica, e (2) as firmas têm maior probabilidade de ajustar o preço quando há grandes choques em seus mercados ou inflação sustentada. No entanto, provou-se difícil introduzir esta forma de ajuste de preços em modelos macroeconômicos completos. Caplin e Leahy (1991) indicaram que as consequências poderiam ser importantes, mas também que muitas simplificações eram necessárias para caracterizar o equilíbrio competitivamente imperfeito com ajuste de preços oneroso, incluindo restrições extremas sobre as regras do banco central, sobre o comportamento dos consumidores e sobre a natureza da demanda por moeda. Assim, embora a precificação dependente do estado seja natural, os modelos existentes têm sido mal adaptados para análise empírica ou exame de regras de política monetária alternativas. Por esse motivo, a ênfase na literatura Nova-Keynesiana tem sido nas regras de ajuste de preços dependentes do tempo (\textit{time-dependent}), que especificam que as firmas têm oportunidades exógenas para o ajuste de preços.

\subsubsection{5.3.1 Uma Abordagem Intertemporal para a Fixação de Preços} Seguindo Calvo (1983), consideramos como uma firma racional selecionaria seu preço hoje, dado que terá que mantê-lo fixo por um intervalo de duração estocástica. Para operacionalizar essa ideia, usamos a notação e estrutura de um estudo recente sobre precificação dependente do tempo e do estado. Como no modelo de competição imperfeita acima, podemos postular um grande número de firmas — tecnicamente um contínuo de firmas — e supor que uma fração $\omega_j$ ajustou seu preço pela última vez há $j$ períodos, para $j = 0, 1, \dots, J-1$. Consequentemente, a probabilidade condicional da data $t$ do próximo ajuste na data $t+j$ é $\omega_{j, t+j}/\omega_{0,t}$. Quando a elasticidade da demanda é assumida constante, de modo que $y_t(z) = [P_t(z)/P_t]^{-\varepsilon} d_t$, com $P_t$ sendo o índice de preços perfeito e $d_t$ um construto de demanda agregada, então o preço ótimo é restrito por:

\begin{equation}
P_t^* = \frac{\varepsilon}{\varepsilon-1} \frac{E_t \sum_{j=0}^{J-1} \beta^j (\Lambda_{t+j}/\Lambda_t) \omega_{j, t+j} (\Psi_{t+j} P_{t+j}^{\varepsilon} d_{t+j})}{E_t \sum_{j=0}^{J-1} \beta^j (\Lambda_{t+j}/\Lambda_t) \omega_{j, t+j} (P_{t+j}^{\varepsilon} d_{t+j})}, \tag{5.9}
\end{equation}

onde $\Psi_{t+j}$ é o custo marginal nominal em $t+j$ e $\beta^j \Lambda_{t+j}/\Lambda_t$ é o fator de desconto para fluxos de caixa contingentes na data $t+j$. A regra geral de ajuste de preços (5.9) deriva de uma equiparação da receita marginal e do custo marginal em um cenário dinâmico e tem uma forma aproximada conveniente que usamos abaixo. Em particular, quando a taxa de inflação está próxima de zero, então $\log P_t^*$ é aproximadamente $\log(\varepsilon/(\varepsilon-1)) + [1/\sum_{h=0}^{J-1} \beta^h \omega_h] \sum_{j=0}^{J-1} \beta^j \omega_j E_t \log \Psi_{t+j}$. Isto é, o preço é uma liderança distribuída descontada do custo marginal nominal esperado, com os pesos relacionados à distribuição de frequência das datas de ajuste de preços. Equivalentemente, denotando o custo marginal real como $\psi_t$ e usando três identidades ($\log \Psi_t = \log P_t + \log \psi_t$, $\log \psi_t = -\log \mu_t$, e $\log \mu = \log(\varepsilon/(\varepsilon-1))$), podemos expressar o preço ótimo como

\begin{equation}
\log P_t^* \approx \frac{1}{\sum_{h=0}^{J-1} \beta^h \omega_h} [ \sum_{j=0}^{J-1} \beta^j \omega_j (E_t \log P_{t+j} + \log (\psi_{t+j}/\psi)) ], \tag{5.10}
\end{equation}

ou seja, como dependendo da trajetória futura do nível de preços e do desvio do custo marginal real de seu nível de estado estacionário.

\subsubsection{5.3.2 O Nível de Preços} Para completar o modelo dinâmico de precificação, precisamos de uma equação que agregue os preços entre as firmas no nível geral de preços. Com todas as firmas que se ajustam na data $t$ escolhendo $P_t^*$, o agregador de preços perfeito é

\begin{equation}
P_t = \left( \sum_{j=0}^{J-1} \omega_{jt} (P_{t-j}^*)^{1-\varepsilon} \right)^{1/(1-\varepsilon)}, \tag{5.11}
\end{equation}

de modo que o nível de preços depende das decisões de precificação e dos padrões de ajuste. Se a variação nos padrões de ajuste for pequena ao longo do ciclo de negócios — como em modelos dependentes do tempo ou alguns modelos dependentes do estado — e a taxa de inflação for baixa, então existe uma aproximação comparável,

\begin{equation}
\log P_t \approx \sum_{j=0}^{J-1} \omega_j \log P_{t-j}^*, \tag{5.12}
\end{equation}

que podemos emparelhar com (5.10). Essas duas equações (5.10) e (5.12) são uma representação conveniente do "bloco de preços" central dos modelos NNS que descrevemos na próxima seção.

\subsubsection{5.3.3 Comparação com o Sistema Dinâmico de Taylor} Baseado na otimização intertemporal e em três simplificações (baixa inflação, elasticidade de demanda constante e pequenas variações nos padrões de ajuste), obtivemos um par de equações log-lineares (5.10) e (5.12) descrevendo a dinâmica de preços. Elas se assemelham amplamente às equações de fixação de salários voltadas para o futuro e de nível de preços voltadas para o passado usadas por Taylor, mas com flexibilidade adicional nos mecanismos de liderança e defasagem distribuídas devido ao uso de um modelo de ajuste estocástico.

Há, no entanto, uma omissão notável: não existem choques de preços em nosso par de expressões comportamentais. Este é um resultado comum na modelagem econômica: a teoria da otimização leva a ver os choques como surgindo de eventos mais primitivos que afetam os tomadores de decisão econômicos. Como veremos, nossa abordagem de otimização permite muitos tipos de eventos que são tipicamente descritos como choques de preços [como, por exemplo, as variações de preços de commodities incluídas por Gordon (1982) em sua especificação empírica]. No entanto, estes exercem uma influência sobre os preços através do custo marginal, em vez de diretamente, de acordo com a teoria desenvolvida na próxima seção.

\section{6. A Nova Síntese: Descrição e Mecânica}

A Nova Síntese Neoclássica é definida por dois elementos centrais. Construindo sobre a nova macroeconomia clássica e a análise RBC, ela incorpora a otimização intertemporal e as expectativas racionais em modelos macroeconômicos dinâmicos. Construindo sobre a economia Nova-Keynesiana, ela incorpora a competição imperfeita e o ajuste de preços oneroso. Como o programa RBC, ela busca desenvolver modelos quantitativos de flutuações econômicas.

A NNS é atualmente exibida em três escalas de modelagem distintas. Primeiro, existem pequenos modelos analíticos que podem ser usados para estudar uma gama de questões teóricas e empíricas, mantendo a tratabilidade suficiente para que possam ser resolvidos à mão. Segundo, existem modelos macroeconômicos de média escala análogos aos desenvolvidos por pesquisadores RBC que estão sendo usados para abordar uma ampla gama de questões positivas e normativas. Terceiro, existe o novo modelo de larga escala FRB/US da economia americana desenvolvido nos últimos anos, que é agora o principal modelo empregado para avaliação de política pelo Conselho do Federal Reserve.

Chamamos o novo estilo de pesquisa macroeconômica de Nova Síntese Neoclássica porque ela herda o espírito da antiga síntese discutida na Seção 2. Os modelos NNS oferecem conselhos de política baseados na ideia de que a rigidez de preços implica que a demanda agregada é um determinante fundamental da atividade econômica real no curto prazo. Os modelos NNS implicam que a política monetária exerce uma influência poderosa sobre a atividade real. Isso tem implicações tanto positivas quanto normativas. De um ponto de vista positivo, a conclusão central é que as flutuações econômicas não podem ser interpretadas ou compreendidas independentemente da política monetária. Isso é verdade não obstante o fato de que o modelo RBC no cerne da NNS atribui um papel potencialmente grande à produtividade, à política fiscal ou a choques de preços relativos. De uma perspectiva normativa, a NNS diz que a demanda agregada deve ser gerida pela política monetária a fim de entregar resultados macroeconômicos eficientes. Em outras palavras, a NNS cria uma demanda urgente por conselhos de política monetária.

A Nova Síntese Neoclássica também fornece esses conselhos. A combinação de fixação de preços racional voltada para o futuro, competição monopolística e componentes RBC na NNS fornece orientação para a política monetária baseada no seguinte raciocínio. Primeiro de tudo, uma política monetária estacionária deve respeitar os determinantes RBC da atividade econômica real em média ao longo do tempo. Isto é, mesmo que o produto possa ser determinado pela demanda período a período na NNS, o produto deve ser determinado pela oferta em média. Segundo, a NNS localiza a transmissão da política monetária para a atividade real em sua influência sobre a razão entre o preço da firma média e o custo marginal de produção, a qual chamamos de \textit{markup} médio. Uma ação de política monetária que aumenta a demanda agregada eleva o custo marginal e reduz o \textit{markup} médio. Esse \textit{markup} médio mais baixo sustenta o aumento no produto e no emprego, porque funciona como uma redução de impostos em um cenário RBC. Terceiro, existe pouco \textit{trade-off} de longo prazo entre inflação e atividade real em baixas taxas de inflação. Ilustrando este ponto, mostramos dentro de uma versão de Taylor de precificação ótima — uma na qual a firma típica ajusta seu preço uma vez por ano — que o imposto de \textit{markup} de estado estacionário é minimizado por uma política monetária que persegue uma inflação próxima de zero. Assim, a recomendação é que a política monetária deve estabilizar a trajetória do nível de preços a fim de manter o produto em seu potencial. Esta política é "ativista" no sentido de que a autoridade deve gerir a demanda agregada para acomodar quaisquer distúrbios do lado da oferta ao produto.

O poder da nova síntese reside na complementaridade de seus componentes Nova-Keynesiano e RBC, que são compatíveis devido à sua dependência compartilhada da microeconomia. A Nova Síntese permite que o conhecimento ganho em estudos Nova-Keynesianos e RBC seja trazido para questões de ciclo de negócios e política monetária em um único modelo coerente. Ao fazê-lo, a nova síntese fortalece nossa compreensão das flutuações econômicas. Esta e as seções subsequentes elaboram sobre as principais características e implicações dos modelos NNS. O restante desta seção cobre algumas preliminares — a mecânica básica dos \textit{markups}, o \textit{markup} médio como um imposto sobre a atividade econômica, os preços relativos como choques de produtividade e o poder e limitações da política monetária. Os princípios NNS e as diretrizes práticas para a política monetária são desenvolvidos nas Seções 7 e 8, respectivamente.

\subsection{6.1 Como a Política Monetária Afeta a Economia Real}

Na nova síntese, a política monetária tem efeitos que se assemelham aos de choques de produtividade e fiscais, produzindo efeitos substituição e riqueza na economia como nos modelos RBC. Variações no \textit{markup} médio cobrado pelas firmas afetam os retornos marginais aos fatores de uma forma que é semelhante aos choques de produtividade ou mudanças nos impostos abrangentes; mudanças nos preços relativos entre as firmas funcionam como os efeitos de nível dos choques de produtividade ou mudanças nos gastos governamentais.

\subsubsection{6.1.1 Markups Marginal e Médio} Duas medidas do \textit{markup} desempenham um papel importante nos modelos da NNS. Como sugerido acima, o \textit{markup} médio do preço sobre o custo marginal desempenha um papel proeminente na transmissão da política monetária. Em qualquer momento, porém, apenas um subconjunto de firmas está ajustando os preços e estabelecendo um novo nível de \textit{markup}, ao qual chamamos de \textit{markup} marginal.

Formalmente, o \textit{markup} marginal é a razão entre o preço e o custo marginal para as firmas que estão ajustando seu preço no período $t$, ou seja,

\begin{equation}
\mu_t^* = \frac{P_t^*}{\Psi_t}. \tag{6.1}
\end{equation}

Sabemos pela Seção 5 que $P_t^*$ depende das expectativas que as firmas em ajuste têm sobre as condições econômicas futuras, incluindo o nível de preços e o custo marginal. O \textit{markup} médio é a razão entre o preço e o custo marginal para a firma média na economia (a razão entre o nível de preços e o custo marginal),

\begin{equation}
\mu_t = \frac{P_t}{\Psi_t}. \tag{6.2}
\end{equation}

Para analisar a determinação da atividade econômica real dentro do período $t$, é o \textit{markup} médio que é central. Deste ponto de vista, é importante enfatizar que o \textit{markup} médio é apenas o recíproco do custo marginal real,

\begin{equation}
\psi_t = \frac{\Psi_t}{P_t} = \frac{1}{\mu_t}. \tag{6.3}
\end{equation}

Assim, a variação pró-cíclica do custo marginal real — que muitos economistas consideram realista — implica diretamente um \textit{markup} médio contracíclico.

Os \textit{markups} médio e marginal podem se mover de forma muito diferente um do outro em resposta a choques. Em resposta a aumentos sustentados na demanda agregada nominal, por exemplo, o \textit{markup} cai para as firmas que não estão ajustando o preço, mas a inflação mais alta motiva as firmas em ajuste a escolher um \textit{markup} mais alto. Assim, o efeito de curto prazo de um aumento sustentado na demanda é que o \textit{markup} marginal sobe e o \textit{markup} médio cai. Retornaremos a este ponto na Seção 7.

\subsubsection{6.1.2 O Markup Médio como um Imposto Distorcivo} As firmas produzem produto com serviços de capital e trabalho. Como são competitivamente monopolísticas, suas demandas por fatores são baseadas na minimização de custos a um nível de produto determinado pela demanda. Uma condição necessária para a minimização de custos é que o valor do produto marginal de cada fator seja igualado ao seu preço de aluguel. Usando $\Psi_t$ para denotar o custo marginal nominal como acima e deixando $W_t$ ser a taxa de salário nominal, a condição de eficiência para o trabalho é $W_t = \Psi_t \partial F(n_t, k_t)/\partial n_t$, e há uma condição comparável para os serviços de capital. Dividindo cada lado desta expressão pelo nível de preços, o salário real é igualado ao custo marginal real vezes o produto marginal do trabalho:

\begin{equation}
w_t = \psi_t \frac{\partial F(n_t, k_t)}{\partial n_t} = \frac{1}{\mu_t} \frac{\partial F(n_t, k_t)}{\partial n_t}, \tag{6.4}
\end{equation}

onde a última igualdade segue diretamente do fato de que o \textit{markup} médio e o custo marginal real são recíprocos. Novamente, uma igualdade semelhante dos preços reais dos fatores e dos produtos marginais de valor real vale para os serviços de capital.

Assim, variações no \textit{markup} médio funcionam exatamente como um imposto abrangente que uma firma deve pagar sobre os insumos de fatores. No caso da demanda por trabalho, por exemplo, o \textit{markup} médio cria uma cunha (\textit{wedge}) entre o salário real e o produto marginal do trabalho, assim como a cunha tributária fez em (4.3). Um \textit{markup} mais alto eleva o imposto implícito sobre o trabalho e o capital.

\subsubsection{6.1.3 Dispersão de Preços Relativos como um Deslocamento de Produtividade} Além do \textit{markup} médio, há uma segunda fonte importante de distorções inerentes aos modelos NNS. Como alguns preços individuais são rígidos, mudanças no nível geral de preços provocam mudanças nos preços relativos. Essa dispersão de preços relativos resulta em uma má alocação do produto agregado entre usos alternativos de bens finais. Para expor essa má alocação, definimos o produto agregado como a soma simples

\begin{equation}
y_t = \int_0^1 y_t(z) dz.
\end{equation}

Suponha ainda, como em (5.5) acima, que a demanda é dada pela especificação de elasticidade constante, $y_t(z) = [P_t(z)/P_t]^{-\varepsilon} d_t$, com $d_t$ sendo o nível de benefício derivado no uso final (consumo ou investimento). Então a distribuição de preços relativos influencia a extensão do benefício de uso final do produto final:

\begin{equation}
d_t = \frac{y_t}{\int_0^1 [P_t(z)/P_t]^{-\varepsilon} dz}.
\end{equation}

As consequências normativas das variações nesta medida composta de preços relativos são análogas às de um choque de produtividade que aumenta o fator total.

\subsection{6.2 O Mecanismo de Transmissão}

A Nova Síntese Neoclássica fornece duas formas complementares de pensar sobre a transmissão das ações de política monetária para a atividade econômica real, que vemos como as abordagens de demanda agregada e de imposto de \textit{markup}.

\subsubsection{6.2.1 Demanda Agregada} De uma perspectiva tradicional, mudanças na quantidade de moeda alteram a demanda agregada, o que exige mudanças na oferta agregada. Quando os modelos NNS são interpretados dessa maneira — considerando a demanda agregada real como determinada pela política monetária — os resultados são sensatos em cada ponto do tempo. No entanto, essa interpretação é incompleta por duas razões: o nível de preços pode responder à política monetária dentro do período, e o foco é deslocado do custo marginal real, que é um elemento importante dos modelos NNS.

\subsubsection{6.2.2 O Imposto de Markup} Uma visão alternativa do mecanismo de transmissão monetária é sugerida pela ideia de que o \textit{markup} pode ser interpretado como um imposto e, em particular, como uma mudança em um imposto geral sobre o produto (vendas) que afeta as recompensas ao capital e ao trabalho. De uma perspectiva RBC, a influência da política monetária sobre a atividade econômica pode ser analisada usando os efeitos relativamente bem compreendidos de mudanças de impostos abrangentes sobre a atividade macroeconômica, que revisamos na Seção 4 acima. Esta visão coloca os movimentos no \textit{markup} médio e no custo marginal real no centro do mecanismo pelo qual a política monetária influencia a atividade econômica real. É igualmente incompleta, no entanto, na medida em que não incorpora a influência do nível de preços sobre o \textit{markup} médio, nem reconhece o papel do custo marginal real na evolução dos preços. No entanto, o \textit{markup} médio continua sendo uma estatística de resumo útil para a transmissão monetária.

\subsection{6.3 O Poder e as Limitações da Política Monetária}

Como sua antecessora homônima, a Nova Síntese Neoclássica vê a política monetária como tendo o potencial de exercer uma grande influência sobre a atividade econômica, embora dentro de limites claramente definidos. Além disso, essa influência pode igualmente ser entendida como operando via distorções, embora diferentes daquelas identificadas na síntese original da década de 1960.

\subsubsection{6.3.1 O Que a Política Monetária Pode Fazer} Para ilustrar o poder da política monetária, é útil estudar o modelo de fixação de preços mais simples possível, um com fixação de preços escalonada em dois períodos. Neste cenário, pode-se supor que a política monetária tenha poder limitado para influenciar a atividade real porque as decisões de precificação são tomadas apenas com um período de antecedência, mas veremos que a política monetária ainda é muito poderosa. No cenário de dois períodos com $\omega_0 = \omega_1 = \frac{1}{2}$, a equação aproximada para o nível de preços (5.12) é $\log P_t = \frac{1}{2}(\log P_t^* + \log P_{t-1}^*)$. A equação de fixação de preços voltada para o futuro (5.10) é

\begin{equation}
\log P_t^* = \frac{1}{1 + \beta} [ \log P_t + \log (\psi_t / \psi) + \beta E_t P_{t+1} + \beta \log E_t (\psi_{t+1} / \psi) ].
\end{equation}

Combinando as equações neste bloco de preços de dois períodos, podemos expressar o nível de preços como

\begin{equation}
\log P_t = \log P_{t-1}^* + \log (\psi_t/\psi) + 2 \sum_{j=1}^{\infty} \beta^j \log E_t (\psi_{t+j}/\psi). \tag{6.5}
\end{equation}

Esta solução de expectativas racionais para o nível de preços exibe duas características importantes que se estendem a modelos de precificação de horizonte mais longo. Primeiro, o nível de preços é parcialmente predeterminado por preços estabelecidos no passado. Segundo, os preços estabelecidos pelas firmas que estão se ajustando atualmente dependem do custo marginal real atual e futuro. De fato, o nível de preços depende de uma liderança distribuída infinita do custo marginal real esperado, embora cada firma deva manter seu preço fixo por apenas dois períodos. As expectativas do custo marginal real futuro importam para a precificação atual porque cada firma sabe que manterá seu preço fixo por algum período de tempo. Além disso, cada firma se preocupa com quais serão os preços no próximo período ao definir seu preço hoje, e por isso se preocupa com quais preços as firmas definirão no próximo período, e assim por diante para o futuro.

A fim de pensar sobre a evolução do nível de preços e do produto neste modelo NNS simples, precisamos entender o comportamento do custo marginal real. Para isso, recorde mais uma vez que o custo marginal real é apenas o recíproco do \textit{markup} médio, então podemos escrever $\log (\psi_t/\psi) = -\log (\mu_t/\mu)$. A análise RBC acima indicou que as variações no imposto de \textit{markup} podem exercer um poderoso efeito inverso sobre o emprego e o produto. Tais efeitos podem ser complexos em um cenário RBC totalmente dinâmico, mas para fins heurísticos, considere a simples relação inversa

\begin{equation}
\log \frac{\mu_t}{\mu} = -\varphi (\log y_t - \log \bar{y}_t), \tag{6.6}
\end{equation}

onde $\bar{y}_t$ é o nível de produto com preços flexíveis, ou seja, aquele obtido em um modelo RBC não competitivo com um \textit{markup} constante $\mu$. Como o custo marginal real é, por sua vez, dado por $\log(\psi_t/\psi) = \varphi(\log y_t - \log \bar{y}_t)$, o parâmetro $\varphi$ é a elasticidade do custo marginal real em relação a um "hiato do produto" (\textit{output gap}).

Agora suponha que o equilíbrio monetário seja dado por uma equação quantitativa como (3.1): $\log M_t = \log P_t + \log y_t - \log v_t$, onde $v_t$ é o processo de velocidade. Substituindo o nível de preços e o produto na equação quantitativa com as expressões do nível de preços e do \textit{markup}, chegamos a uma expressão relacionando o estoque de moeda aos \textit{markups} atuais e futuros esperados:

\begin{equation}
\log (M_t) = \log(P_{t-1}^*) + \log \bar{y}_t + \log (v_t) - \frac{1+\varphi}{\varphi} \log (\mu_t/\mu) - 2 \sum_{j=1}^{\infty} \beta^j \log E_t (\mu_{t+j}/\mu). \tag{6.7}
\end{equation}

Assim, desde que a autoridade monetária siga uma política que sustente uma liderança distribuída determinada de \textit{markups} esperados (uma condição relativamente fraca), a expressão precedente indica que ela pode escolher o estoque de moeda para produzir um padrão arbitrário de pequenas variações no \textit{markup} médio ao longo do tempo. A autoridade monetária teria uma alavancagem semelhante sobre a trajetória do \textit{markup}, do custo marginal real e do produto em modelos NNS mais gerais também.

Uma forma de resumir esse poder é que a autoridade monetária pode escolher um processo estocástico estacionário arbitrário para o imposto de \textit{markup} em relação a uma média $\mu$. No entanto, a autoridade monetária pode produzir variações no \textit{markup} médio apenas aceitando as implicações para os preços e para a moeda. Em particular, a estabilização do \textit{markup} e a estabilização do nível de preços estão intimamente relacionadas nos modelos NNS, um ponto ao qual retornaremos quando discutirmos o papel da política monetária.

\subsubsection{6.3.2 O Que a Política Monetária Não Pode Fazer} No entanto, a analogia com a tributação é incompleta. Embora a autoridade monetária possa escolher como o imposto de \textit{markup} se move ao longo do tempo, há pouco que ela possa fazer para afetar o nível de estado estacionário do \textit{markup}, porque a NNS incorpora a fixação de preços voltada para o futuro. Como discutimos na próxima seção, em baixas taxas de inflação, o nível do \textit{markup} de estado estacionário é quase invariante à taxa de inflação e, portanto, é essencialmente determinado pela extensão do poder de monopólio no setor privado. Além disso, existem algumas restrições entre o curto e o longo prazo, como em qualquer modelo de expectativas racionais. Quanto mais persistentes forem os movimentos planejados da autoridade monetária no imposto de \textit{markup}, maiores serão suas consequências inflacionárias.

\section{7. Princípios Norteadores para a Política Monetária}

A Nova Síntese Neoclássica faz a forte recomendação de que um banco central deve buscar uma inflação próxima de zero. Nesta seção, explicitamos os princípios subjacentes a essa prescrição. Para concretude e simplicidade, trabalhamos dentro do modelo de fixação de preços dependente do tempo desenvolvido acima. Os princípios são suficientemente gerais, no entanto, para que orientem a política monetária em outros modelos NNS também. O papel da política monetária na nova síntese deriva de duas fontes. Primeiro, a estrutura microeconômica subjacente sugere que é desejável estabilizar o \textit{markup} médio, evitando uma fonte de distorções variáveis no tempo para a macroeconomia. Segundo, o comportamento de fixação de preços voltado para o futuro torna viável o desenho de políticas simples que realizarão essa estabilização.

\subsection{7.1 A Taxa Ótima de Inflação}

Quais são as implicações da nova síntese para a taxa ótima de inflação? Embora uma análise completa deste tópico esteja além do escopo do presente estudo, podemos identificar várias características fundamentais que são importantes. Primeiro, a taxa de inflação afeta a distribuição de preços relativos em qualquer modelo com rigidez de preços, o que por sua vez tem efeitos sobre o valor de uso final do produto que descrevemos acima. Estes são minimizados quando há inflação zero. Segundo, o \textit{markup} médio depende da taxa de inflação: no exemplo que estudamos mais detalhadamente abaixo, o \textit{markup} médio é minimizado a uma taxa de inflação que é próxima de zero. Terceiro, se recursos são gastos no ajuste de preços, então estes são minimizados na inflação zero. Assim, nestes três fundamentos, a incorporação da competição imperfeita e da rigidez de preços leva à sugestão de que uma taxa de inflação próxima de zero é desejável.

No entanto, Friedman (1969) argumentou anteriormente que era desejável ter deflação esperada, de modo que a taxa de juros nominal de curto prazo fosse zero. Assim, uma análise completa da taxa ótima de inflação deve equilibrar os benefícios monetários da desinflação com os custos de distorção associados à deflação.

\subsubsection{7.1.1 Efeitos sobre os Markups} As primeiras análises keynesianas reconheciam que a inflação estável corroeria o poder de mercado das firmas, sugerindo benefícios para a inflação sustentada. No entanto, modelos dinâmicos de fixação de preços sugerem, na melhor das hipóteses, uma pequena taxa de inflação positiva nestes fundamentos. Além disso, esses modelos de fixação de preços também sugerem que taxas maiores de inflação aumentarão, em vez de diminuir, os \textit{markups} médios devido aos efeitos da inflação esperada. Usamos a Figura 3 para exibir os principais ingredientes desta conclusão. Primeiro, em qualquer modelo com preços rígidos, a inflação positiva corrói mecanicamente os preços relativos das firmas que não estão se ajustando ou, equivalentemente, haverá preços relativos mais altos para aquelas firmas que estão se ajustando. Para fornecer uma ideia da importância quantitativa deste canal, o painel A mostra o efeito da inflação sobre $P/P^*$ com uma elasticidade de demanda de 11 e o escalonamento de 4 trimestres do ajuste de preços sugerido por alguns trabalhos de Taylor ($\omega_j = 0,25$ para $j=0, 1, 2, 3$). Uma taxa de inflação anual de 10\% reduz $P/P^*$ em cerca de 4\%. Segundo, confrontada com uma situação de maior inflação de estado estacionário, uma firma racional fixadora de preços tem um incentivo para elevar seu \textit{markup} marginal. Usando os mesmos valores de parâmetros acima, o painel B mostra que uma taxa de inflação de 10\% faz com que $\mu^* = P^*/\Psi$ aumente para 1,15 a partir do nível de inflação zero $\varepsilon/(\varepsilon-1) = 1,10$. Assim, as firmas elevam o \textit{markup} marginal substancialmente em resposta à inflação antecipada.

O \textit{markup} médio incorpora tanto os efeitos de erosão inflacionária quanto de inflação esperada, uma vez que é simplesmente o produto $\mu = (\mu^*) (P/P^*)$. Consequentemente, o painel C exibe o efeito combinado da inflação sobre o \textit{markup} médio, rendendo três resultados de interesse. Primeiro, o menor valor do \textit{markup} médio ocorre a uma taxa de inflação positiva, mas esta taxa não é muito diferente de zero. Segundo, o efeito da inflação sobre o \textit{markup} é positivo em taxas de inflação mais altas. Terceiro, o efeito geral da inflação sobre o \textit{markup} médio é quantitativamente muito pequeno próximo da inflação zero. No entanto, inflações maiores na verdade elevam em vez de diminuir o \textit{markup} médio: aumentar a taxa de inflação para 10\% ao ano a partir de zero produz um aumento em $\mu$ de $\varepsilon/(\varepsilon-1) = 1,10$ para 1,1044. Assim, com pequenas inflações ou deflações, a autoridade monetária não pode influenciar muito o \textit{markup} médio no modelo NNS.

\subsubsection{7.1.2 Distorções de Preços Relativos pela Inflação} Se não houver inflação no estado estacionário, então todos os preços relativos serão um e o valor de uso final do produto será maximizado. Além disso, pequenas mudanças nos preços relativos próximas deste ponto inicial não terão efeito sobre a razão $d_t/y_t$, de modo que não haverá efeito de produtividade de pequenos ciclos de negócios ou pequenas taxas de inflação ou deflação. No entanto, usando uma elasticidade de demanda de $\varepsilon=11$ e escalonamento de preços em 4 trimestres como acima, calculamos que uma taxa de inflação anual de 10\% reduzirá o valor de uso final do produto em 0,4\% e, mais geralmente, exibirá a relação entre inflação e distorções de preços relativos no painel D da Figura 3. Assim, o arcabouço NNS indica um custo social direto quantitativamente importante da inflação sustentada decorrente de distorções de preços relativos.

Tomando esses achados concernentes ao \textit{markup} médio e ao tamanho das distorções de preços relativos juntamente com a observação da Seção 4 de que existem ganhos relativamente pequenos ao reduzir a inflação de zero para a regra de Friedman, o modelo NNS recomenda que a autoridade monetária busque uma taxa de inflação próxima de zero. Uma vez que os efeitos de produtividade das distorções de preços relativos são menores próximo da inflação zero, no que segue focamos apenas nos movimentos no \textit{markup} médio ao considerar a resposta da macroeconomia a vários choques.

\subsection{7.2 Política Monetária e o Ciclo de Negócios}

Como a política monetária deve variar ao longo do ciclo de negócios? Argumentamos que o objetivo da autoridade monetária deve ser produzir uma trajetória constante para o \textit{markup} médio ou, equivalentemente, para o custo marginal real. Embora a constância do \textit{markup} seja um objetivo \textit{ad hoc}, ela é atraente para nós por três razões relacionadas. Primeiro, ela traz a mesma resposta da economia real a vários choques que surgiria se todos os preços fossem perfeitamente flexíveis. Segundo, ela corresponde à suavização tributária conforme recomendado na literatura de finanças públicas. Finalmente, é consistente com o foco tradicional sugerido da política monetária, que é eliminar hiatos entre o produto atual e um nível de produto potencial (capacidade) variável no tempo.

Nossa recomendação equivale a uma \textit{política monetária neutra} no sentido de que ela mantém o \textit{markup} médio em seu nível de estado estacionário e faz o modelo NNS se comportar como uma economia RBC não competitiva. A política monetária neutra acomoda choques que alterariam os níveis de equilíbrio de produto e emprego com preços flexíveis, tais como mudanças na produtividade, política fiscal e mudanças nos preços relativos internacionais, e alguns que não o fariam, tais como mudanças na demanda por moeda.

\subsubsection{7.2.1 Política Monetária Neutra e Dinâmica de Preços} A dinâmica de preços NNS envolve componentes voltados para o futuro e para o passado, conforme discutido na seção anterior. Em uma primeira aproximação [como em (5.10)], uma firma em ajuste define seu preço em

\begin{equation}
\log P_t^* = \frac{1}{\sum_{h=0}^{J-1} \beta^h \omega_h} \left[ \sum_{j=0}^{J-1} \beta^j \omega_j (E_t \log P_{t+j} + \log (\psi_{t+j}/\psi)) \right]
\end{equation}

Em uma primeira aproximação [como em (5.12)], o nível de preços é $\log P_t = \sum_{j=0}^{J-1} \omega_j \log P_{t-j}^*$. Sob o requisito de política monetária neutra, o custo marginal real é constante agora e em todas as datas futuras, de modo que a fixação de preços depende apenas da trajetória futura esperada do nível de preços. Consequentemente, as duas equações dinâmicas implicam uma equação de diferença expectacional que pode ser resolvida para determinar as implicações de nível de preços e inflação da política monetária neutra.

\subsubsection{7.2.2 A Desejabilidade de um Nível de Preços Constante} O resultado de referência é que um nível de preços constante é uma política monetária neutra. Isto é, se definirmos $\log P_t^* = \log P_t = \log \bar{P}$ em todas as datas nas equações de preços (5.10) e (5.12), então um valor presente do custo marginal real deve ser esperado como sendo zero em todas as datas, o que só pode ser satisfeito por um nível constante de custo marginal real. Existem duas formas equivalentes de declarar essa conclusão. Diretamente, uma autoridade monetária comprometida em buscar uma trajetória de crescimento para o nível de preços deve fazê-lo mantendo o custo marginal real constante ou, equivalentemente, um \textit{markup} médio constante. Alternativamente, pode-se dizer que uma autoridade monetária comprometida com a política neutra deve buscar uma taxa de inflação constante.

\subsubsection{7.2.3 O Processo de Oferta de Moeda Sustentando Resultados Neutros} Sob uma política monetária neutra, o produto comporta-se de acordo com um ciclo de negócios real competitivamente monopolístico com um \textit{markup} constante $\mu$ diante de choques de tecnologia, política fiscal e preços relativos internacionais. A política neutra elimina os hiatos de produto, tornando $y_t = \bar{y}_t$ em todas as datas. Sob uma política neutra, a autoridade monetária acomoda variações na demanda por moeda para garantir que excessos ou escassez de moeda não criem distúrbios de demanda agregada. Para trabalhar as implicações para a oferta de moeda, suponha que a trajetória do nível de preços sob política neutra seja dada por $\log P_t = \log P_{t-1} + \pi$, onde $\pi$ é a taxa de inflação tendencial. Como a inflação é constante, as variações na taxa de juros real ($r$) e nominal ($R$) são idênticas ($R_t = r_t + \pi$). Então, se a demanda por moeda for $\log M_t = \log P_t + m_y \log y_t - m_R R_t + v_t$, o estoque de moeda deve ser

\begin{equation}
\log M_t = \log P_t + m_y \log \bar{y}_t - m_R (r_t + \pi) + v_t. \tag{7.1}
\end{equation}

Isto é: a autoridade monetária deve acomodar movimentos no produto e nas taxas de juros obtidos no modelo RBC, e choques de velocidade também.

\subsubsection{7.2.4 Condições Iniciais e Transições de Inflação} As equações de precificação ótima permitem prontamente uma caracterização da política monetária neutra sob condições mais gerais. Duas décadas atrás, Edmund Phelps e Guillermo Calvo estudaram o problema da desinflação em um modelo básico de salário fixo com uma estrutura matemática semelhante a (5.10) e (5.12). Duas características fundamentais da política monetária neutra estendem-se à economia da desinflação. Primeiro, o \textit{markup} médio deve ser constante ao longo do tempo, o que equivale a exigir que a decisão de ajuste de preços dependa apenas da trajetória futura esperada do nível de preços. Segundo, a trajetória do nível de preços é apenas uma função das decisões de ajuste de preços tomadas em várias datas. Quando resolvemos a equação de diferença expectacional resultante assumindo que a taxa de inflação de estado estacionário é zero, a solução "estável" é da forma

\begin{equation}
\log P_t^* - \log P_{t-1}^* = \sum_{j=1}^{J-2} \sigma_j (\log P_{t-j}^* - \log P_{t-j-1}^*), \tag{7.2}
\end{equation}

onde os coeficientes $\sigma_j$ são funções dos parâmetros $\omega_j$ e $\beta$. Ou seja, existe uma trajetória única de ajustes de preços que deve ocorrer se houver um \textit{markup} médio constante. Existem várias implicações desta fórmula de desinflação neutra de Phelps-Calvo. Primeiro, a política monetária neutra poderia equivalentemente ser declarada como uma regra para a taxa de crescimento dos preços recém-definidos, $\pi_t^* = \log P_t^* - \log P_{t-1}^*$. Segundo, dado que determinamos a taxa de crescimento $\pi_t^*$ necessária para uma política monetária neutra, podemos usar a equação do nível de preços (5.12) para determinar a trajetória de transição neutra para a taxa de inflação medida, $\pi_t = \log P_t - \log P_{t-1} = \sum_{j=0}^{J-1} \omega_j \pi_{t-j}^*$. No modelo de Taylor de 4 trimestres, uma transição neutra para $\pi = 0$ assume a forma

\begin{equation}
\pi_t^* = -0,43 \pi_{t-1}^* - 0,12 \pi_{t-2}^*. \tag{7.3}
\end{equation}

Isto é, quando começamos em um estado estacionário inflacionário com uma taxa trimestral de inflação de, digamos, 2,5\% (para que a taxa de inflação anual seja inicialmente de 10\%), então deve haver uma queda de preços por parte das firmas em ajuste igual a $-(0,42 + 0,12) \times 0,025 = -0,01375$ no período de impacto de uma desinflação neutra. Esta queda de preços é necessária para estabilizar o \textit{markup} médio dado os aumentos de preços passados incorporados no sistema. Com esta ação de política agressiva, a taxa de inflação real $\pi_t = \frac{1}{4} (\pi_t^* + \pi_{t-1}^* + \pi_{t-2}^* + \pi_{t-3}^*)$ cai de 2,5\% para cerca de 1,5\% no período de impacto da política e subsequentemente declina rapidamente para zero ao longo de um ano.

\subsubsection{7.2.5 Controle Imperfeito do Nível de Preços} Também podemos operacionalizar a política monetária neutra quando a autoridade monetária tem controle imperfeito do nível de preços. Em tal cenário, a autoridade monetária não pode alcançar o controle perfeito do imposto de \textit{markup}, mas pode evitar que ele varie em valor esperado. Ou seja, sua regra de política pode fazer $E_{t-1} \sum_{j=0}^{\infty} \beta^j \omega_j \log(\mu_{t+j}/\mu) = 0$. Os resultados precedentes aplicam-se então ao componente esperado da política monetária, com um choque de ajuste de preço adicional introduzido na análise. Ou seja, com controle imperfeito do nível de preços, a política monetária neutra assume a forma

\begin{equation}
\pi_t^* = \sum_{j=1}^{J-2} \sigma_j \pi_{t-j}^* + \xi_t, \tag{7.4}
\end{equation}

onde $\xi_t = \log P_t^* - E_{t-1} \log P_t^*$. Assim, o banco central acomoda alguma porção dos erros de meta do nível de preços, como na análise de Taylor.

\subsubsection{7.2.6 Comparação de Metas de Inflação e Regras de Nível de Preços} Muitos bancos centrais perseguem metas de inflação que permitem o desvio da base (\textit{base drift}) no nível de preços. Em nosso cenário, um retorno a uma trajetória de nível de preços fixo é indesejável, uma vez que requer variações no \textit{markup} médio. Podemos usar a análise precedente para quantificar quanto desvio da base é desejável no cenário com fixação de preços escalonada em 4 trimestres dada em (7.3). Suponha que a informação incompleta leve a um erro de meta, $\xi_t > 0$, do qual a autoridade monetária toma conhecimento ao final do período atual. No cenário atual (7.3), o efeito desejável de longo prazo sobre o nível de preços é simplesmente

\begin{equation}
\frac{\xi_t}{1 - \sum \sigma_j} = \frac{\xi_t}{1 - (-0,43 - 0,12)} \approx 0,6 \xi_t.
\end{equation}

Assim, a autoridade monetária permite que cerca de seis décimos de um erro de previsão de preços se transmitam para o nível geral de preços no longo prazo.

\section{8. A Prática da Política Monetária}

Embora a estabilidade de preços tenha sido sugerida há muito tempo como um objetivo primário para a política monetária, uma série de questões surgiu sobre sua desejabilidade e viabilidade prática. Esta seção aborda quatro preocupações principais usando a abordagem da Nova Síntese Neoclássica. Primeiro, os efeitos do petróleo e outros choques de preços de commodities. Segundo, Milton Friedman e outros monetaristas questionaram a desejabilidade das metas de inflação com base em sua leitura da história monetária e nas defasagens longas e variáveis. Terceiro, Novos Keynesianos como John Taylor sugeriram a existência de \textit{trade-offs} importantes entre a variabilidade do produto e da inflação. Quarto, os banqueiros centrais rotineiramente se preocupam com as táticas de uso de seu instrumento de política preferido, uma taxa de juros de curto prazo.

\subsection{8.1 Um Choque do Petróleo no Modelo da Nova Síntese}

Os choques do petróleo representam um problema difícil para a política monetária porque podem criar inflação e desemprego ao mesmo tempo. No entanto, esse problema torna os choques do petróleo particularmente instrutivos para ilustrar a mecânica do arcabouço NNS e seu poder prescritivo para a política monetária. A análise também destaca a complementaridade do raciocínio RBC e keynesiano que é herdada pelos modelos NNS.

É natural pensar em um choque do petróleo como uma restrição na oferta de petróleo disponível para uso na produção de bens finais. As firmas produzem produto combinando (após os custos fixos) serviços de capital e trabalho com petróleo. Uma vez que as firmas são competitivamente monopolísticas, o produto é determinado pela demanda. Para qualquer nível de demanda final, as demandas ótimas por fatores exigem que o custo marginal de produção de produto para uma firma seja o mesmo para um aumento em qualquer um dos três fatores de produção. Por analogia a (6.4), o uso ótimo de energia exige que

\begin{equation}
\Psi_t a_t \frac{\partial F}{\partial q_{et}} = P_{et},
\end{equation}

onde $q_{et}$ é a quantidade de insumo de energia (petróleo) e $P_{et}$ é o preço nominal do petróleo. Uma firma escolhe os usos ótimos de fatores tomando os preços dos fatores como dados. Em equilíbrio geral, os preços dos fatores se ajustam para equilibrar os mercados de fatores e, ao influenciar o \textit{markup}, os ajustes nos preços dos fatores também ajudam a equilibrar o mercado de bens finais.

Uma vez que o nível de preços é rígido, o produto é governado pela demanda agregada no curto prazo. Assim, precisamos assumir uma posição sobre como a demanda agregada se comportará para dizer como o sistema responde ao choque do petróleo. Para fins ilustrativos, nossa estratégia é perguntar o que a política de demanda agregada deve fazer, e assumir que a política monetária sustenta esse nível de demanda agregada.

Balizamos a resposta de política ótima com o raciocínio RBC. Por construção, o modelo RBC competitivo padrão responderia eficientemente ao choque do petróleo. Uma condição necessária para que o modelo NNS responda de forma eficiente é que ele também mantenha um \textit{markup} constante. Assim, a NNS recomenda que a política monetária deve visar estabilizar o \textit{markup} contra o choque do petróleo, não acomodando qualquer parte do aumento do preço do petróleo em uma inflação mais alta.

Com a política monetária neutra em vigor, podemos perguntar como o modelo NNS responderia ao choque do petróleo. Nos níveis iniciais de insumos de fatores, produto e preço, o aumento no preço do petróleo eleva o custo marginal nominal e, portanto, corta o \textit{markup}. Para que a política monetária restaure o \textit{markup} ao seu nível inicial, a política deve deprimir a demanda agregada e cortar o emprego. Da perspectiva keynesiana, tal recomendação soa como acrescentar insulto à injúria — fazendo o emprego cair justamente quando os custos dos materiais estão altos. No entanto, o raciocínio RBC diz que a economia deve produzir menos quando o custo marginal de produção está temporariamente alto. Esse raciocínio também sugere que a extensão do corte adequado na demanda depende da persistência esperada do choque do petróleo. Um choque que se espera ser temporário tem pouco efeito riqueza sobre a oferta de trabalho e a demanda de consumo. Ele aumenta principalmente o custo de oportunidade do trabalho atual em relação ao lazer e do lazer atual em relação ao lazer futuro. Assim, a política monetária deve agir para cortar a demanda agregada temporariamente para refletir esses custos de oportunidade. As taxas de juros reais devem subir para sustentar a política monetária neutra.

Um choque do petróleo que se espera ser altamente persistente, por outro lado, agiria como um choque de produtividade negativo persistente, criando um grande efeito riqueza negativo que compensaria o efeito substituição sobre a oferta de trabalho. Relativamente pouco declínio no emprego poderia ser exigido neste caso. Mas seria apropriado que a política monetária provocasse um corte no consumo proporcional ao declínio na produtividade devido à falta de disponibilidade de petróleo.

Para resumir, pode-se razoavelmente perguntar por que, na prática, os choques do petróleo têm sido inflacionários. Primeiro, na medida em que os produtos de petróleo são produzidos em mercados competitivos e comprados diretamente pelos consumidores, o aumento no preço do petróleo entra diretamente no nível de preços sem ser intermediado pelas firmas produtoras de bens no setor de preços rígidos da economia. Estabilizar o nível de preços contra esses choques de preços diretos exigiria perseguir uma política de demanda agregada restritiva o suficiente para empurrar a demanda e o emprego para baixo no setor de preços rígidos, aumentando assim o \textit{markup} ali. O raciocínio NNS não recomenda aumentar o \textit{markup} no setor de preços rígidos para estabilizar o nível geral de preços. A política deve ser acomodatícia a tais choques de preços diretos, especialmente porque são choques de preços relativos cujo efeito sobre a inflação é temporário. Segundo, e igualmente importante, os bancos centrais podem relutar em deixar as taxas de juros reais subirem acentuadamente, especialmente quando um choque de custos está prejudicando a economia. As consequências inflacionárias dos choques dos preços do petróleo foram provavelmente exacerbadas pelas tentativas dos bancos centrais de suavizar as taxas de juros nominais com um crescimento monetário excessivamente expansionista.

\subsection{8.2 As Metas de Inflação são Práticas?}

Os economistas monetários há muito pensam que a estabilidade de preços tem muito a recomendá-la como o objetivo primário para a política monetária, e recentemente uma série de bancos centrais adotou metas de inflação explícitas como um guia para a política. Tem sido menos claro, no entanto, que as metas de inflação poderiam desempenhar um papel útil como um objetivo de política imediato e um critério de desempenho. Usando a NNS, revisamos argumentos práticos que foram avançados contra as metas de inflação por Friedman (1960) e outros. Argumentamos que essas objeções são indevidamente pessimistas quando se reconhece o papel dos preços rígidos e da credibilidade do banco central na fixação de preços.

\subsubsection{8.2.1 Interpretando a Experiência Histórica} A visão de Friedman baseia-se em grande parte em seu trabalho sobre a história monetária dos Estados Unidos com Anna Schwartz, no qual encontraram defasagens no efeito da política monetária como sendo longas e variáveis, variando entre meio ano a mais de dois anos. Raciocinando com base nos dados históricos, Friedman observou que "o nível de preços... poderia ser um guia eficaz apenas se fosse possível prever, primeiro, os efeitos não monetários sobre o nível de preços por um período considerável de tempo no futuro, e segundo, o tempo que levará em cada instância particular para as ações monetárias terem seu efeito..." Ele concluiu que "... a ligação entre mudanças de preços e mudanças monetárias em períodos curtos é muito frouxa e imperfeitamente conhecida para tornar a estabilidade do nível de preços um objetivo e guia razoavelmente inequívoco para a política".

A inferência de Friedman sobre a conveniência das metas de inflação parece pessimista demais. Em primeiro lugar, nenhum dos dados da história monetária dos EUA foi extraído de um regime de política guiado pela busca proposital da estabilidade de preços. O padrão-ouro anterior à Primeira Guerra Mundial era aquele em que a inflação tendencial era pequena pelos padrões de hoje. Mas os Estados Unidos não tinham banco central, e o crescimento monetário era fortemente influenciado por pânicos bancários em várias ocasiões, e por fluxos de ouro governados pelo balanço de pagamentos e pelo acaso de novas descobertas e técnicas de mineração. Como consequência, a variabilidade do nível de preços no curto prazo foi bastante significativa em alguns momentos durante o período.

Após a fundação do Federal Reserve, houve inflação durante a Primeira Guerra Mundial seguida por uma deflação acentuada após a guerra; então os preços se estabilizaram na década de 1920, e o nível de preços caiu cerca de um terço de 1929 a 1933. A inflação da Segunda Guerra Mundial não foi revertida subsequentemente, e em vez disso a nação entrou em um período no qual o nível de preços mais que triplicou nas três décadas seguintes à Guerra da Coreia. Os modelos NNS implicam que as ligações entre preços e produto dependem sensivelmente do regime monetário. Como a história monetária dos EUA tem sido uma sucessão de regimes monetários muito diferentes, a NNS preveria exatamente o tipo de instabilidade aparente no efeito da moeda encontrado por Friedman e Schwartz. As descobertas de Robert Gordon de equações de preços empíricas radicalmente diferentes em diferentes períodos de amostra são uma manifestação do mesmo tipo de instabilidade dependente do regime.

\subsubsection{8.2.2 O Papel da Credibilidade} Se as inferências dos dados históricos podem ser enganosas, podemos fazer algumas conjecturas sobre as metas de baixa inflação no modelo NNS com base no papel da credibilidade do banco central no processo de fixação de preços. De acordo com (5.10), por exemplo, a fixação de preços onerosa implica que as firmas se preocupam com uma liderança distribuída do nível de preços e do custo marginal real ao definir o preço de hoje. Quando um regime de metas de inflação é perfeitamente crível, lideranças distribuídas fixas tanto de preços quanto de custo marginal real (o recíproco do \textit{markup}) ancoram o comportamento atual de fixação de preços. Adicione a isso algum escalonamento da fixação de preços, e a presunção é que a credibilidade para a baixa inflação tende a ser auto-reforçadora em grande medida, porque em tal ambiente as firmas pensarão menos sobre a inflação e estarão menos nervosas com ela. Esta confiança seria reforçada ainda mais por um mandato legislativo tornando a baixa inflação uma prioridade para a política monetária.

A questão principal para um banco central comprometido com a baixa inflação é quão "indulgentes" os fixadores de preços provavelmente serão com os erros de política. Erros ocorrerão inevitavelmente devido à informação imperfeita sobre a economia. Mas tais erros teriam pouco efeito se corrigidos a tempo, precisamente devido à lentidão na fixação de preços. É claro que um banco central que permitisse que os erros se acumulassem por algum motivo, de modo que a inflação começasse a subir significativamente, poderia transformar a liderança distribuída na equação de preços de uma âncora estabilizadora em uma fonte de sustos inflacionários desestabilizadores.

Os sustos inflacionários (\textit{inflation scares}) são fáceis de entender sob a perspectiva da nova síntese. Um banco central tem um incentivo para trapacear em seu compromisso com a estabilidade de preços no modelo NNS porque uma ação de política monetária pode reduzir a distorção do \textit{markup} e aumentar o emprego. Chari, Kehoe e Prescott (1989), por exemplo, poderiam argumentar que um banco central sem uma tecnologia de pré-compromisso não poderia sustentar um equilíbrio de baixa inflação de forma alguma. No mínimo, seu argumento sugere que o incentivo para trapacear torna os fixadores de preços hipersensíveis a erros de política de uma forma que torna um equilíbrio de baixa inflação extremamente frágil.

Parece-nos que o raciocínio NNS, juntamente com desenvolvimentos recentes na política monetária, enfraquece consideravelmente a força de tal ponto. Pensamos que bancos centrais como o Federal Reserve hoje internalizam amplamente os custos de longo prazo de trapacear. Como resultado do Fed de Volcker ter assumido a responsabilidade pela inflação no final da década de 1970 e ter tido sucesso em reduzi-la, o Fed é agora amplamente considerado responsável pela inflação. Além disso, a experiência de baixa inflação desde então demonstrou claramente os benefícios de longo prazo da estabilidade de preços. Portanto, acreditamos que a tentação para o Fed trapacear em seu compromisso de baixa inflação é muito mais fraca do que no passado.

\subsection{8.3 Inflação e Variabilidade do Produto}

Embora seu modelo de contratos escalonados e sobrepostos não exiba nenhum \textit{trade-off} de longo prazo no nível de inflação e no nível de produto, Taylor (1980) mostrou que ele implica um \textit{trade-off} entre a variância do produto e a variância da inflação. Com base nisso, Taylor argumentou que os ciclos de negócios podem ser reduzidos apenas aceitando uma maior variabilidade da inflação.

Uma vez que os modelos NNS incorporam o tipo de comportamento de fixação de preços assumido por Taylor, surge a questão de saber se tais modelos também apresentam aos formuladores de políticas uma escolha difícil entre a variabilidade da inflação e do produto. A questão é de interesse mais do que acadêmico, uma vez que recai sobre uma das questões mais importantes no banco central hoje: o desenho de um mandato legislativo para a política monetária. A maioria dos especialistas concorda que alguma forma de mandato claro melhoraria a eficácia da política ao fixar as expectativas de inflação e aumentar a responsabilidade do banco central. A nova síntese apoia tal raciocínio. Mas não há acordo sobre se existe um \textit{trade-off} ou, se existir, sobre como permitir isso em um mandato.

\subsubsection{8.3.1 Existe um Trade-off?} Recorde nosso princípio de que a política monetária nos modelos NNS deve visar manter o \textit{markup} constante no nível baixo consistente com uma inflação próxima de zero. Assim, a política monetária deve compensar os choques na demanda agregada. Tais ações de política não apenas manteriam o produto no potencial, mas estabilizariam os preços também. Por outro lado, a política monetária deve acomodar choques de produtividade, levando em conta quaisquer efeitos associados na oferta de trabalho e no estoque de capital. Caso contrário, um hiato de produto se abriria, o que faria o \textit{markup} variar. Não há \textit{trade-off} em nenhum destes casos — a política deve estabilizar tanto o \textit{markup} quanto os preços em resposta a choques de demanda ou de produtividade. Mesmo para um choque do petróleo, a sociedade claramente não enfrenta nenhum \textit{trade-off} se o petróleo for um insumo intermediário. Vimos acima que o melhor resultado é manter a estabilidade de preços e reduzir a demanda em resposta ao declínio na produtividade.

O que aconteceria em um modelo NNS com um setor de produção de bens de preços flexíveis ao lado do setor competitivo monopolístico de preços rígidos, no qual os choques poderiam impactar a inflação diretamente? Claramente, tal modificação não alteraria a conclusão com respeito aos choques de demanda agregada ou de produtividade, uma vez que estes ainda deveriam ser compensados ou acomodados, respectivamente. A flexibilidade de preços adicionada, no entanto, complica a resposta a um choque do petróleo, porque a restrição na oferta de petróleo faz com que o preço do petróleo suba em relação a outros preços. Se a política devesse deprimir a demanda agregada apenas o suficiente para manter os preços estáveis no setor de preços rígidos, os preços dos produtos intensivos em petróleo no setor de preços flexíveis ainda subiriam. O banco central poderia reduzir a demanda agregada o suficiente para evitar que o nível geral de preços (preços flexíveis mais rígidos) subisse, mas então ele elevaria o \textit{markup} e criaria um hiato de produto no setor de preços rígidos.

Assim, a política pareceria enfrentar um \textit{trade-off} entre inflação e variabilidade do produto em relação a choques de preços relativos. Mas mesmo aqui, o raciocínio NNS fornece uma saída. Praticamente falando, a nova síntese sugere que um banco central deve visar estabilizar um índice apenas de preços rígidos, um índice de preços "núcleo" (\textit{core}). Esta visão concorda bem com a ênfase keynesiana em um índice de núcleo em vez de um índice de custo de vida geral, e a recomendação monetarista de estabilizar um índice de longo prazo e ignorar tais movimentos de preços relativos como os choques dos preços do petróleo. Quando definimos a medida de preços que um banco central deve estabilizar como um índice de núcleo de preços rígidos, encontramos mais uma vez que não há \textit{trade-off} de política entre inflação e variabilidade do produto.

\subsection{8.4 Implementação Tática de Política}

A nova síntese sugere que um banco central deve perseguir uma política ativista para visar a inflação. Existem grandes dificuldades em implementar uma regra de política ativista, muitas delas bem conhecidas e debatidas há muito tempo entre economistas monetários e banqueiros centrais, algumas das quais foram abordadas acima. Nosso propósito nesta seção é fazer alguns pontos adicionais sugeridos pela nova síntese para pensar sobre a implementação prática da política.

\subsubsection{8.4.1 Política de Taxa de Juros} Os bancos centrais invariavelmente usam uma taxa de juros de curto prazo como seu instrumento de política monetária. A nova síntese diz que os banqueiros centrais devem gerir um regime de metas de baixa inflação fazendo com que a taxa nominal de curto prazo mimetize a taxa real de curto prazo que seria gerada por um modelo RBC bem especificado com um \textit{markup} baixo e constante. O raciocínio RBC é indispensável para pensar sobre quanto e em que direção a taxa real deve ser movida em resposta a um choque. Por exemplo, até mesmo a direção da resposta apropriada da taxa real a um choque de produtividade depende da duração esperada do choque, como vimos acima quando discutimos o choque do petróleo.

Como outro exemplo do valor do raciocínio RBC, considere isto. Recentemente, uma possível aceleração no crescimento da produtividade foi citada como uma razão pela qual o Federal Reserve não precisa elevar as taxas de juros reais de curto prazo para manter a baixa inflação. De fato, o componente RBC padrão do modelo NNS sugere, no mínimo, que as taxas reais teriam que subir um para um com um aumento no crescimento da produtividade tendencial, por exemplo, um aumento de 50 pontos-base na taxa de crescimento seria acompanhado por um aumento de 50 pontos-base nas taxas de juros reais. É importante notar que as taxas teriam que subir mesmo se a economia estivesse de outra forma operando em um nível potencial de PIB não inflacionário.

\section{9. Resumo e Conclusões}

Os modelos da Nova Síntese Neoclássica são complexos, pois envolvem otimização intertemporal, expectativas racionais, competição monopolística, ajuste de preços oneroso e precificação dinâmica, além de um papel importante para a política monetária. Nossos objetivos principais no artigo foram triplos: motivar os componentes separados da nova síntese, apresentar um arcabouço conceitual para pensar sobre os modelos NNS e usar esse aparato para desenvolver recomendações para a política monetária.

Dois \textit{insights} fundamentais estão no cerne do nosso arcabouço. Primeiro, os mecanismos keynesianos e RBC podem ser vistos como operando através de canais um pouco diferentes. Mantendo o \textit{markup} médio constante, a mecânica do modelo NNS assemelha-se à de um modelo RBC puro, embora não competitivo. Por outro lado, a influência keynesiana da demanda agregada sobre o emprego e o produto funciona encolhendo ou aumentando o \textit{markup}, que age como um imposto distorcivo sobre a atividade econômica.

Segundo, o ajuste dinâmico de preços oneroso significa que as firmas ajustam o preço de acordo com uma liderança distribuída esperada do nível de preços e do custo marginal real, onde o nível de preços é uma média dos preços atuais e daqueles estabelecidos no passado. Mostramos que a equação de precificação voltada para o futuro e uma expressão de nível de preços formam um bloco de preços que pode ser resolvido para expressar a taxa de inflação como uma função de preços estabelecidos no passado, custo marginal real atual e uma liderança distribuída do custo marginal real esperado. Como o custo marginal real é o inverso do \textit{markup}, a evolução da inflação no modelo NNS depende de forma importante dos \textit{markups} atuais e futuros esperados.

A política monetária neutra recomendada na nova síntese segue diretamente dos \textit{insights} acima e da ideia de que o \textit{markup} deve ser mantido constante. A constância do \textit{markup} é atraente porque entrega a mesma resposta da economia real a vários choques que surgiria se todos os preços fossem perfeitamente flexíveis. Mostramos que o \textit{markup} de estado estacionário deve ser minimizado a uma taxa de inflação próxima de zero, e argumentamos que a maioria dos benefícios para a troca monetária seria realizada na inflação próxima de zero também. Assim, descobrimos que a meta de inflação próxima de zero era tanto desejável quanto viável no modelo NNS.

Mesmo que a nova síntese herde muito do espírito da antiga, ela difere agudamente em termos do papel da política monetária. Economistas trabalhando dentro da síntese da década de 1960 eram pessimistas quanto a domar a inflação, vendo a inflação como tendo um ímpeto próprio e flutuando com mudanças incontroláveis na psicologia dos fixadores de preços. A nova síntese também vê as expectativas como críticas para o processo inflacionário, mas vê as expectativas como passíveis de gestão por uma regra de política monetária.

A nova síntese tem muito a dizer sobre a implementação prática das metas de inflação. Como as expectativas de \textit{markups} futuros desempenham um papel fundamental no processo gerador de inflação, uma meta de inflação bem-sucedida exige um compromisso crível com a baixa inflação, para que as expectativas de constância do \textit{markup} ancorem a equação geradora de inflação. A fim de manter a constância do \textit{markup}, a política monetária deve acomodar movimentos no PIB potencial provocados por forças RBC, tais como produtividade, política fiscal ou choques de custos de materiais. A acomodação deve ser bidimensional. Primeiro, o crescimento da moeda deve satisfazer os movimentos induzidos na demanda por moeda. Segundo, a autoridade monetária deve mover seu instrumento de taxa de juros nominal de curto prazo para apoiar os movimentos da taxa de juros real de curto prazo exigidos pelas forças RBC subjacentes. Ironicamente, apesar do fato de que os efeitos keynesianos da política monetária sobre a atividade real são poderosos nos modelos NNS, a política monetária é melhor quando elimina inteiramente os efeitos keynesianos.

Os pesquisadores apenas arranharam a superfície ao pensar sobre os modelos NNS: tais modelos certamente serão elaborados e melhorados no futuro. Olhando para trás: os modelos NNS devem melhorar nossa compreensão dos resultados macroeconômicos durante períodos inflacionários voláteis, tais como aquele que se estende de meados da década de 1960 até o início da década de 1980, quando tanto grandes choques de política monetária quanto grandes choques de oferta foram importantes. Olhando para frente: à medida que os Estados Unidos e outros países ao redor do mundo mantêm a inflação baixa, as forças do lado da oferta devem assumir uma importância tão grande quanto as forças do lado da demanda para o ciclo de negócios. Esperamos que os modelos NNS se tornem cada vez mais importantes no fornecimento de conselhos de política monetária em tal ambiente.

\end{document}